% ═══════════════════════════════════════════════════════════════════════════════
%   MODULO SYNTHESIS — El Cálculo de Kuratowski
%   Single-column textbook format · pdfLaTeX
% ═══════════════════════════════════════════════════════════════════════════════
\documentclass[11pt, a4paper, openany]{book}

% ─── Geometry ─────────────────────────────────────────────────────────────────
\usepackage[
  top=2.5cm, bottom=2.5cm,
  left=3cm, right=3cm,
  headheight=14pt
]{geometry}

% ─── Fonts ────────────────────────────────────────────────────────────────────
\usepackage{fontspec}
\defaultfontfeatures{Ligatures=TeX}

% ─── Language ────────────────────────────────────────────────────────────────
\usepackage{polyglossia}
\setmainlanguage{spanish}

% ─── Colors ───────────────────────────────────────────────────────────────────
\usepackage{xcolor}
\definecolor{structure}{HTML}{1A535C}   % Deep Teal
\definecolor{main}{HTML}{C6892A}        % Golden Ochre
\definecolor{second}{HTML}{2E4057}      % Slate Blue
\definecolor{third}{HTML}{8B2635}       % Burgundy
\definecolor{fourth}{HTML}{3A7D44}      % Forest Green
\definecolor{mainlight}{HTML}{FDF6E3}   % Solarized Base3
\definecolor{secondlight}{HTML}{EDF2F4}
\definecolor{thirdlight}{HTML}{F9EBEC}
\definecolor{fourthlight}{HTML}{EBF5EC}

% ─── Packages ─────────────────────────────────────────────────────────────────
\usepackage{amsmath,amssymb,amsthm,mathtools}
\usepackage{graphicx,booktabs,enumitem,microtype,float}
\usepackage{fancyhdr,titlesec,titletoc}
\usepackage[numbers,sort&compress]{natbib}
\bibliographystyle{unsrtnat}
\usepackage{tcolorbox}
\tcbuselibrary{skins,breakable,theorems}
\usepackage{etoolbox}
\usepackage{hyperref}
\hypersetup{
  colorlinks=true,
  linkcolor=structure, citecolor=second, urlcolor=second,
  bookmarksnumbered=true, pdfstartview=FitH
}

% ─── Macros ───────────────────────────────────────────────────────────────────
\newcommand{\lP}{\ell_{\mathrm{P}}}
\newcommand{\dS}{d_{\mathrm{S}}}
\newcommand{\MPl}{M_{\mathrm{Pl}}}
\newcommand{\stoneref}[1]{%
  \smallskip\noindent
  {\small\color{structure!70}\textit{$\blacktriangleright$ Lectura
  l\'ogica (Ap\'endice~\ref{app:stone-duality}):} #1}%
  \smallskip}
\renewcommand{\chaptername}{Modulo}

% ─── Chapter/Section Styling ─────────────────────────────────────────────────
\titleformat{\chapter}[display]
  {\normalfont\Large\bfseries\color{structure}}
  {\small\color{structure!60}\MakeUppercase{\chaptertitlename}\ \thechapter}
  {4pt}{\Huge}
\titlespacing*{\chapter}{0pt}{-20pt}{20pt}

\titleformat{\section}
  {\normalfont\Large\bfseries\color{structure}}{\thesection}{0.6em}{}
\titlespacing*{\section}{0pt}{12pt}{6pt}

\titleformat{\subsection}
  {\normalfont\large\bfseries\color{structure!80!black}}{\thesubsection}{0.5em}{}
\titlespacing*{\subsection}{0pt}{8pt}{4pt}

% ─── Headers ──────────────────────────────────────────────────────────────────
\pagestyle{fancy}
\fancyhf{}
\fancyhead[LE,RO]{\small\thepage}
\fancyhead[LO]{\small\nouppercase{\rightmark}}
\fancyhead[RE]{\small\nouppercase{\leftmark}}
\renewcommand{\headrulewidth}{0.3pt}

% ─── Theorem Environments ─────────────────────────────────────────────────────
\newcounter{thm}[chapter]
\renewcommand{\thethm}{\thechapter.\arabic{thm}}

% Definition Box
\newtcolorbox[use counter=thm]{definition}[1][]{%
  enhanced,breakable,
  colback=mainlight,colframe=main,
  colbacktitle=main,coltitle=white,
  fonttitle=\small\bfseries,
  title={Definici\'on~\thethm\ifx&#1&\else: #1\fi},
  arc=1.5pt,boxrule=0.6pt,
  left=8pt,right=8pt,top=6pt,bottom=6pt,
  before skip=10pt,after skip=10pt
}

% Theorem Box
\newtcolorbox[use counter=thm]{theorem}[1][]{%
  enhanced,breakable,
  colback=secondlight,colframe=second,
  colbacktitle=second,coltitle=white,
  fonttitle=\small\bfseries,
  title={Teorema~\thethm\ifx&#1&\else: #1\fi},
  arc=1.5pt,boxrule=0.6pt,
  left=8pt,right=8pt,top=6pt,bottom=6pt,
  before skip=10pt,after skip=10pt
}

% Axiom Box
\newtcolorbox[use counter=thm]{axiom}[1][]{%
  enhanced,breakable,
  colback=thirdlight,colframe=third,
  colbacktitle=third,coltitle=white,
  fonttitle=\small\bfseries,
  title={Axioma~\thethm\ifx&#1&\else: #1\fi},
  arc=1.5pt,boxrule=0.6pt,
  left=8pt,right=8pt,top=6pt,bottom=6pt,
  before skip=10pt,after skip=10pt
}

% Result Box
\newtcolorbox{result}[1][]{%
  enhanced,breakable,
  colback=fourthlight,colframe=fourth,
  fontupper=\itshape,
  arc=1.5pt,boxrule=0.6pt,
  left=8pt,right=8pt,top=6pt,bottom=6pt,
  before skip=10pt,after skip=10pt,
  title={Resultado Clave},
  coltitle=fourth!50!black, fonttitle=\small\bfseries,
  attach boxed title to top left={yshift=-2mm,xshift=2mm},
  boxed title style={colback=fourthlight,frame hidden},
  #1
}

% Introduction/Overview Box
\newenvironment{introduction}[1][Sinopsis del M\'odulo]{%
  \begin{tcolorbox}[
    enhanced,colback=structure!5,colframe=structure,
    colbacktitle=structure,coltitle=white,
    fonttitle=\small\bfseries,title={#1},
    arc=2pt,boxrule=0.7pt,
    left=8pt,right=8pt,top=6pt,bottom=6pt,
    before skip=12pt,after skip=12pt
  ]\itshape\begin{itemize}[leftmargin=*,itemsep=2pt,parsep=0pt]
}{%
  \end{itemize}\end{tcolorbox}
}

% Proof environment
\renewenvironment{proof}[1][Demostraci\'on]{%
  \par\noindent\textit{#1.}\quad
}{\hfill$\square$\par\medskip}

\setlist{nosep,leftmargin=*}
\setlist[itemize,1]{label={\color{structure}\textbullet}}
\setlist[enumerate,1]{label={\color{structure}\arabic*.}}

% ─── Title Page Macro ────────────────────────────────────────────────────────
\makeatletter
\newcommand{\subtitle}[1]{\gdef\@subtitle{#1}}
\newcommand{\extrainfo}[1]{\gdef\@extrainfo{#1}}
\gdef\@subtitle{}\gdef\@extrainfo{}
\renewcommand{\maketitle}{%
  \begin{titlepage}\centering
    \vspace*{4cm}
    {\Huge\bfseries\color{structure}\@title\par}\vspace{0.6cm}
    \ifx\@subtitle\empty\else{\Large\color{gray!80}\@subtitle\par}\fi
    \vspace{2cm}
    {\color{main}\rule{0.6\textwidth}{2pt}\par}\vspace{2cm}
    {\LARGE\@author\par}\vspace{1cm}
    {\large\@date\par}\vfill
    \ifx\@extrainfo\empty\else
      \begin{minipage}{0.8\textwidth}\centering\itshape\color{gray!60}\@extrainfo\end{minipage}\vspace{3cm}
    \fi
  \end{titlepage}
}
\makeatother

% ═══════════════════════════════════════════════════════════════════════════════
\title{Modulo Synthesis}
\subtitle{El C\'alculo de Kuratowski}
\author{Juan Pablo Silva Alvarado}
\date{\today}
\extrainfo{``El c\'alculo continuo es la pintura; el C\'alculo de Kuratowski es el lienzo.''}

\begin{document}
\maketitle
\frontmatter
\tableofcontents
\mainmatter

% ███████████████████████████████████████████████████████████████████████████████
%   PROLOGUE
% ███████████████████████████████████████████████████████████████████████████████

\chapter*{Pr\'ologo: La Anatom\'ia del Continuo}
\addcontentsline{toc}{chapter}{Pr\'ologo: La Anatom\'ia del Continuo}
\markboth{Pr\'ologo}{Pr\'ologo}

Durante m\'as de tres siglos, la f\'isica ha hablado el idioma de lo
infinitamente suave. Desde que Newton y Leibniz concibieron la derivada~\cite{newton1687},
y Riemann generaliz\'o la geometr\'ia de las superficies
curvas~\cite{riemann1854}, nuestra comprensi\'on del universo se ha
cimentado sobre una suposici\'on maravillosamente \'util: la
divisibilidad infinita. En este paradigma, el
espacio es un lienzo continuo, el tiempo fluye sin interrupciones, y el
c\'alculo opera sobre un l\'imite diferencial ($dx \to 0$) que se
desvanece suavemente en el abismo de lo infinitamente peque\~no.

Esta presunci\'on nos permiti\'o modelar el cosmos. Nos dio las
herramientas para imaginar esferas conc\'entricas de perspectiva que se
extienden hacia el infinito, anidando nuevos universos y
permiti\'endonos concebir la majestuosidad de un espacio que se curva, se
contrae y rebota sobre s\'i mismo. La geometr\'ia de Riemann y el
c\'alculo diferencial no son un error; son, por el contrario, la
aproximaci\'on estad\'istica m\'as exitosa en la historia del
pensamiento humano~\cite{penrose2003}.

Sin embargo, cuando la gravedad de Riemann intenta hablar con la
mec\'anica cu\'antica, la ilusi\'on de la suavidad se
fractura~\cite{dewitt1967,goroff1986}. Los infinitos que antes eran una
herramienta matem\'atica se convierten en singularidades
patol\'ogicas~\cite{hawking1970,penrose1965}. El continuo se rompe.

\bigskip

\textit{Modulo Synthesis: El C\'alculo de Kuratowski} no es un intento de
destruir el c\'alculo cl\'asico, sino de revelar la maquinaria que opera
debajo de \'el. Su prop\'osito es demostrar que las matem\'aticas
continuas que hemos amado y utilizado durante siglos son el l\'imite
macrosc\'opico de una realidad fundamentalmente discreta, finita y
estructurada como un Grafo Ac\'iclico Dirigido
(DAG)~\cite{bombelli1987,finkelstein1969,thooft1979}.

En este libro, no abandonaremos los s\'imbolos que definen a la
f\'isica te\'orica. En cambio, les devolveremos su rigor ontol\'ogico
redefiniendo su significado a la escala de Planck:

\begin{itemize}
  \item \textbf{El Diferencial ($dx$):} Ya no representa un l\'imite
    geom\'etrico que tiende a cero, sino la unidad fundamental
    indivisible del cambio. Un diferencial es un v\'inculo causal
    primario (un salto topol\'ogico de cobertura, denotado como
    $\prec^{\!*}$). Es el ``p\'ixel'' del espaciotiempo causal.

  \item \textbf{La Integral ($\int$):} Integrar sobre un volumen ya no
    es sumar \'areas infinitesimales bajo una curva continua. En el
    C\'alculo de Kuratowski, la integral es una operaci\'on de conteo
    estricto. Es la cardinalidad de los eventos contenidos dentro de un
    intervalo causal (el Diamante de Alexandrov).

  \item \textbf{La M\'etrica ($g_{\mu\nu}$):} La m\'etrica que dicta
    c\'omo se curvan las distancias en la Relatividad General se revela
    como una propiedad emergente. Su valor local est\'a dictado
    directamente por la densidad de los Prismas Causales ($K_{2,N}$).
    Donde hay materia (nudos topol\'ogicos), el n\'umero de caminos se
    multiplica, la informaci\'on se retrasa, y el observador
    macrosc\'opico percibe esto como una curvatura del continuo y la
    presencia de masa inercial.
\end{itemize}

\bigskip

La geometr\'ia de Kuratowski nos ense\~na que el espacio no es el
contenedor de la realidad, sino la red misma de sus relaciones l\'ogicas.
Este volumen es el manual de traducci\'on. Es el puente anal\'itico que
nos permite partir de un solo bit de informaci\'on causal discreta y,
paso a paso, recuperar las infinitas esferas de perspectiva de Riemann.

\begin{center}
\itshape
El c\'alculo continuo es la pintura; el C\'alculo de Kuratowski es el
lienzo. Veamos de qu\'e est\'a hecho.
\end{center}

\stoneref{Las traducciones f\'isica$\,\leftrightarrow\,$l\'ogica
que se insin\'uan a lo largo de este volumen no son analog\'ias
pedag\'ogicas: son \textbf{isomorfismos categ\'oricos} garantizados
por la Dualidad de Stone~\cite{stone1936}. El
Ap\'endice~\ref{app:stone-duality} lo demuestra formalmente.}

\bigskip
\begin{center}
\rule{0.4\textwidth}{0.5pt}
\end{center}

% ███████████████████████████████████████████████████████████████████████████████
%   PART I : LOS ELEMENTOS DEL CÁLCULO DISCRETO
% ███████████████████████████████████████████████████████████████████████████████
\part{Los Elementos del C\'alculo Discreto}

% ═══════════════════════════════════════════════════════════════════════════════
\chapter{El Diferencial Discreto: Los Elementos Base}
% ═══════════════════════════════════════════════════════════════════════════════
% CORRESPONDENCIA CENTRAL:
%   Continuo: dx → 0  (límite infinitesimal)
%   Discreto: ≺*      (vínculo de cobertura indivisible)

\begin{introduction}
  \item El evento causal como \'atomo indivisible de espaciotiempo.
  \item El v\'inculo de cobertura $\prec^{\!*}$ como diferencial discreto.
  \item Longitud y tiempo propio como conteo de v\'inculos.
  \item El l\'imite de Planck como axioma topol\'ogico, no como constante
        arbitraria.
\end{introduction}

El \'exito irrazonable del c\'alculo continuo en la f\'isica cl\'asica
se basa en una ilusi\'on fundamental: la creencia de que el
espaciotiempo es un sustrato infinitamente divisible. En la geometr\'ia
de Riemann, un evento es un punto adimensional definido por coordenadas
reales $(x^\mu)$, y el cambio se mide a trav\'es de un diferencial
$dx^\mu$ en el l\'imite donde $dx \to 0$.

Sin embargo, la mec\'anica cu\'antica y la topolog\'ia subyacente del
universo nos proh\'iben alcanzar ese cero absoluto. En este cap\'itulo
establecemos el primer pilar del C\'alculo de Kuratowski: la
redefinici\'on estricta de los elementos
base~\cite{bombelli1987,sorkin2003,surya2019}.

% ─────────────────────────────────────────────────────────────────────────────
\section{El Evento como Nodo L\'ogico}

En nuestra formulaci\'on, abandonamos la noci\'on de un ``espacio
contenedor''~\cite{malament1977,hawking1976,kronheimer1967}. El
universo no es un escenario donde ocurren las cosas; es la red de las
cosas que ocurren.

\begin{definition}[El Evento Fundamental]
\label{def:event}
Un \textbf{punto} en el C\'alculo de Kuratowski no posee coordenadas
aprior\'isticas. Se define como un nodo~$u$ perteneciente a un
Grafo Ac\'iclico Dirigido (DAG) causalmente
ordenado~\cite{bombelli1987,myrheim1978}
$\mathcal{C} = (V, \prec)$, donde:
\begin{enumerate}
  \item $V$ es un conjunto finito (o localmente finito) de
        \textbf{eventos}.
  \item $\prec$ es un orden parcial estricto: irreflexivo, transitivo y,
        por la aciclicidad del grafo, asim\'etrico.
  \item Para cualesquiera $u, v \in V$, el intervalo
        $\{w \in V : u \prec w \prec v\}$ es finito
        (\textbf{finitud local}).
\end{enumerate}
El ``volumen'' de un evento no es cero sino la unidad fundamental de
informaci\'on de Planck~($\lP^4$).
\end{definition}

\noindent
Esta definici\'on tiene una consecuencia inmediata: las singularidades
de la Relatividad General~\cite{penrose1965,hawking1973}---centros de
agujeros negros, el instante $t = 0$ del Big Bang---dejan de existir
como objetos matem\'aticos v\'alidos, porque la densidad de informaci\'on est\'a topol\'ogicamente
acotada por la cardinalidad de~$V$. No se puede ``comprimir infinitos
eventos en volumen cero'' cuando cada evento ocupa, por axioma, un
volumen irreducible~$\lP^4$.

% ─────────────────────────────────────────────────────────────────────────────
\section{El L\'imite del Diferencial: $dx \to \prec^{\!*}$}

Si el espacio no es liso, la operaci\'on fundamental del c\'alculo
diferencial---la derivada---requiere una nueva base ontol\'ogica. ?`Qu\'e
significa ``moverse una distancia infinitamente peque\~na'' si el espacio
est\'a pixelado?

\begin{definition}[El Diferencial Discreto (V\'inculo de Cobertura)]
\label{def:covering}
El an\'alogo exacto del diferencial de arco~$ds$ (o el intervalo
temporal~$dt$) es el \textbf{v\'inculo de cobertura causal}
$\prec^{\!*}$. Un evento~$u$ \textbf{cubre} a un evento~$v$
($u \prec^{\!*} v$) si y solo si:
\[
  u \prec v
  \quad\text{y}\quad
  \nexists\; w \in V \;\text{tal que}\; u \prec w \prec v.
\]
La relaci\'on de cobertura es el salto causal m\'inimo: el paso
m\'as peque\~no posible en el universo. No existe ``medio v\'inculo''.
\end{definition}

\stoneref{El v\'inculo de cobertura $\prec^{\!*}$ \emph{es} la
implicaci\'on at\'omica en el \'algebra de Heyting dual
(Teorema~\ref{thm:stone-kuratowski}, punto~1). El diferencial
discreto es una unidad indivisible de \emph{razonamiento}: un paso
inferencial que no admite descomposici\'on.}

\noindent
El diagrama de Hasse del conjunto
causal~\cite{davey2002}---la reducci\'on transitiva del DAG---retiene
\'unicamente los v\'inculos de cobertura~$\prec^{\!*}$, descartando
las aristas redundantes de largo alcance. Es la representaci\'on can\'onica del espaciotiempo
discreto: todo lo que la red sabe est\'a codificado en sus v\'inculos
de cobertura.

Por lo tanto, la notaci\'on cl\'asica del l\'imite diferencial:
\begin{equation}
  \lim_{\Delta x \to 0} \frac{\Delta y}{\Delta x}
  \label{eq:classical-limit}
\end{equation}
carece de sentido f\'isico real. En el C\'alculo de Kuratowski, la
derivada no es el c\'alculo de una tangente imaginaria, sino la
medici\'on discreta de c\'omo fluye la informaci\'on a trav\'es de un
\'unico salto topol\'ogico~$\prec^{\!*}$. El ``cero'' del c\'alculo
tradicional es, en realidad, el ``uno'' de la topolog\'ia causal:
\begin{equation}
  dx \;\longrightarrow\; \prec^{\!*}
  \qquad\text{(un v\'inculo de cobertura, indivisible).}
  \label{eq:dx-to-link}
\end{equation}

% ─────────────────────────────────────────────────────────────────────────────
\section{La Longitud y el Tiempo Propio}

En la Relatividad de Einstein, el tiempo propio $\tau$ experimentado
por un observador a lo largo de una l\'inea de mundo se calcula
integrando el diferencial de arco continuo~$ds$. En nuestra
topolog\'ia, la longitud de una curva se reduce a un problema de
conteo estricto.

\begin{definition}[Cadena Causal]
\label{def:chain}
Una \textbf{cadena causal} es una secuencia totalmente ordenada de
eventos conectados por v\'inculos de cobertura:
\[
  \gamma \;=\; (e_0 \prec^{\!*} e_1 \prec^{\!*} \cdots \prec^{\!*} e_n),
  \qquad e_i \in V,\quad 0 \le i \le n.
\]
La \textbf{longitud} de la cadena es el n\'umero de v\'inculos (no de
nodos): $|\gamma| = n$. Denotamos por $\Gamma(A,B)$ el conjunto de
todas las cadenas causales entre dos eventos $A, B \in V$ con
$A \prec B$.
\end{definition}

\begin{definition}[Tiempo Propio Discreto]
\label{def:proper-time}
Sea $\gamma \in \Gamma(A,B)$ una cadena causal. El \textbf{tiempo
propio discreto} a lo largo de~$\gamma$ es:
\begin{equation}
  \tau(\gamma) \;=\; |\gamma|
  \;=\; \bigl|\{(e_i, e_{i+1}) : 0 \le i < n\}\bigr|
  \;=\; n.
  \label{eq:proper-time}
\end{equation}
Cada v\'inculo de cobertura contribuye exactamente \textbf{una
unidad} de duraci\'on causal. No existe fracci\'on de v\'inculo.
\end{definition}

\stoneref{$\tau(\gamma) = |\gamma|$ es la \emph{profundidad de
demostraci\'on}: el n\'umero de pasos inferenciales at\'omicos en la
cadena de razonamiento $\gamma$. El tiempo es la longitud de una
prueba~\cite{prawitz1965}.}

\noindent
La correspondencia con el continuo es directa:
\[
  \tau_{\text{continuo}}
  \;=\; \int_\gamma ds
  \quad\longrightarrow\quad
  \tau_{\text{discreto}}
  \;=\; \sum_{i=0}^{n-1} 1
  \;=\; n.
\]
La legitimidad f\'isica de esta identificaci\'on descansa sobre un
resultado cl\'asico de la teor\'ia de conjuntos causales:

\begin{theorem}[Myrheim--Meyer~\cite{myrheim1978,meyer1988,brightwell1991}]
\label{thm:myrheim-meyer}
Sea $\mathcal{C} = (V, \prec)$ un conjunto causal obtenido por
\emph{sprinkling} de Poisson con densidad~$\rho$ en un espaciotiempo
$d$-dimensional. Para dos eventos $A \prec B$, sea
$\ell(A,B) = \max_{\gamma \in \Gamma(A,B)} |\gamma|$ la longitud de la
cadena m\'axima. Entonces, en el l\'imite de alta densidad:
\begin{equation}
  \frac{\ell(A,B)}{\rho^{1/d}}
  \;\xrightarrow{\;\rho \to \infty\;}
  \; m_d \;\tau(A,B),
  \label{eq:myrheim-meyer}
\end{equation}
donde $\tau(A,B)$ es el tiempo propio geod\'esico del continuo y
$m_d$ es una constante que depende solo de la dimensi\'on~$d$.
\end{theorem}

\noindent
El Teorema~\ref{thm:myrheim-meyer} establece que el conteo de
v\'inculos de cobertura a lo largo de la cadena m\'as larga converge
al tiempo propio riemanniano. La cadena m\'axima es la geod\'esica
(v\'ease el Cap\'itulo~4).

\medskip
Esta identificaci\'on desmitifica la relatividad del tiempo. El tiempo
se dilata para un objeto masivo no porque ``viaje por un espacio
curvo'', sino porque su informaci\'on est\'a atrapada procesando las
bifurcaciones internas de los Prismas Causales~($K_{2,N}$), obligando
a su cadena a consumir v\'inculos de cobertura en estructura local
antes de avanzar hacia el futuro global. La masa no frena al
tiempo; lo \emph{ocupa}.

\begin{theorem}[Dilataci\'on Temporal como Retraso Topol\'ogico]
\label{thm:time-dilation}
Sean $A \prec B$ dos eventos del conjunto causal y sean
$\gamma_{\text{vac}}$ y $\gamma_{\text{mat}}$ dos cadenas m\'aximas
entre~$A$ y~$B$ que atraviesan, respectivamente, una regi\'on de
vac\'io y una regi\'on con densidad de Prismas Causales
$\nu_{\text{mat}} > 0$. Si la presencia de Prismas genera
$\Delta N > 0$ v\'inculos internos adicionales por unidad de
profundidad causal que no contribuyen al avance global, entonces:
\begin{equation}
  |\gamma_{\text{mat}}^{\text{global}}|
  \;<\;
  |\gamma_{\text{vac}}^{\text{global}}|,
  \label{eq:time-dilation}
\end{equation}
donde $|\gamma^{\text{global}}|$ denota el n\'umero de v\'inculos de
cobertura que avanzan hacia el futuro causal com\'un de~$A$ y~$B$
(excluyendo los v\'inculos consumidos en rutas internas de los Prismas).
\end{theorem}

\begin{proof}
El n\'umero total de v\'inculos disponibles para un observador
entre~$A$ y~$B$ est\'a acotado por la longitud de la cadena
m\'axima $\ell(A,B)$, que depende solo de la geometr\'ia global
(Teorema~\ref{thm:myrheim-meyer}). En la regi\'on de vac\'io, todos
los v\'inculos avanzan causalmente:
$|\gamma_{\text{vac}}^{\text{global}}| = |\gamma_{\text{vac}}|$.
En la regi\'on con materia, cada Prisma $K_{2,N_j}$ desv\'ia al
caminante por $N_j$ v\'inculos intermedios que conectan los polos del
Prisma sin avanzar en el orden causal global. As\'i:
\[
  |\gamma_{\text{mat}}^{\text{global}}|
  \;=\; |\gamma_{\text{mat}}|
  \;-\; \sum_j N_j.
\]
Puesto que $\sum_j N_j > 0$, se obtiene
la desigualdad~\eqref{eq:time-dilation}.
\end{proof}

\stoneref{Derivaci\'on independiente v\'ia teor\'ia de la
prueba~\cite{prawitz1965}: las demostraciones que utilizan lemas
(introducci\'on de corte) son m\'as largas que las demostraciones
directas. Los Prismas son lemas; la dilataci\'on temporal es el costo
de la indirecci\'on inferencial.}

% ─────────────────────────────────────────────────────────────────────────────
% TODO — §1.3: El Límite de Planck como Axioma Topológico
% En el continuo, ℓ_P es una constante dimensional derivada de ℏ, G, c
% \cite{planck1899,garay1995}.
% Aquí invertimos la lógica: la discretitud es el axioma primario, y ℓ_P
% es la consecuencia dimensional de la densidad fundamental ρ del sprinkling.
% Relación: ρ = 1/ℓ_P⁴.
% Esto elimina la jerarquía de escalas como input y la convierte en output.
% (Cubierto parcialmente por §1.4 y Axioma~\ref{ax:planck-density}.)

% ─────────────────────────────────────────────────────────────────────────────
\section{Los Par\'ametros de la Existencia}

A diferencia de la f\'isica cl\'asica, donde las constantes fundamentales
se obtienen mediante observaci\'on experimental, el C\'alculo de
Kuratowski postula que los valores que gobiernan la estructura del
espaciotiempo son \textbf{constantes topol\'ogicas} derivadas de la
teor\'ia de grafos. No hay par\'ametros libres que ajustar; cada valor
se sigue de una restricci\'on combinatoria sobre el DAG causal.

\begin{axiom}[L\'imite de Valencia $\mathcal{V}_{\max}$]
\label{ax:valence-limit}
Todo evento $u \in V$ del conjunto causal satisface una cota superior
sobre su grado total en el diagrama de Hasse:
\begin{equation}
  \deg(u) \;=\; d^+(u) + d^-(u) \;\leq\; \mathcal{V}_{\max}.
  \label{eq:valence-limit}
\end{equation}
\end{axiom}

\noindent
El l\'imite de valencia define la \textbf{saturaci\'on del flujo
causal}: el n\'umero m\'aximo de v\'inculos de cobertura que pueden
confluir en un evento. Este par\'ametro es el responsable de que la
dimensi\'on espectral se estabilice en un valor finito ($\dS \to 4$)
en vez de divergir. Si la valencia fuera ilimitada, la probabilidad
de retorno $P(t)$ no decaer\'ia como ley de potencias y la
dimensi\'on espectral no estar\'ia bien definida. El valor
$\mathcal{V}_{\max} = 15$ emerge de la condici\'on de que el
\emph{sprinkling} de Poisson en un diamante causal de 4 dimensiones
produce, en media, $\langle\deg\rangle \approx 2d = 8$ v\'inculos de
cobertura por evento, con una cola poissoniana~\cite{kingman1993} que se trunca
naturalmente a $\sim 2\langle\deg\rangle$.

\begin{axiom}[Umbral de Kuratowski $K_{\text{threat}}$]
\label{ax:kuratowski-threshold}
Un nodo~$w$ constituye una \textbf{amenaza topol\'ogica} para un
Prisma Causal $\mathcal{P}(u, v, W)$ si est\'a conectado
simult\'aneamente a ambos polos y a al menos $K_{\text{threat}}$
intermedios:
\begin{equation}
  w \in V \setminus (W \cup \{u,v\})
  \quad\text{con}\quad
  |\{u, v\} \cap N(w)| = 2
  \;\;\text{y}\;\;
  |W \cap N(w)| \;\geq\; K_{\text{threat}},
  \label{eq:kuratowski-threshold}
\end{equation}
donde $N(w)$ denota los vecinos de~$w$ en el diagrama de Hasse. Cuando
esta condici\'on se cumple, el nodo~$w$ es absorbido por contracci\'on
de v\'ertice en el polo m\'as cercano, preservando la planaridad del
grafo.
\end{axiom}

\noindent
El valor $K_{\text{threat}} = 2$ se deriva directamente del
\textbf{Teorema de Kuratowski}~\cite{kuratowski1930,diestel2017}: un
grafo es planar si y solo si no contiene un menor de $K_5$ ni de
$K_{3,3}$. Un nodo conectado a los 2
polos y a 2 intermedios de un Prisma $K_{2,N}$ forma un subgrafo con
$2 + 2 + 1 = 5$ v\'ertices mutuamente alcanzables---una
configuraci\'on que amenaza con producir un menor de~$K_5$. La
absorci\'on es el mecanismo m\'inimo que neutraliza esta amenaza. Es el
an\'alogo discreto del \textbf{confinamiento de color}: la topolog\'ia
proh\'ibe que un nodo externo ``rompa'' la estructura bipartita del
Prisma, del mismo modo que la QCD proh\'ibe aislar un quark.

\begin{axiom}[Profundidad de Hasse $d_{\max}$]
\label{ax:hasse-depth}
La b\'usqueda de Prismas Causales y amenazas topol\'ogicas opera dentro
de un horizonte local de profundidad~$d_{\max}$ en el diagrama de
Hasse:
\begin{equation}
  d_{\max} \;=\; 4.
  \label{eq:hasse-depth}
\end{equation}
Todo v\'inculo a distancia de Hasse $> d_{\max}$ de un evento dado
queda fuera de su cono de influencia inmediata y se trata como ruido
de fondo.
\end{axiom}

\noindent
La profundidad $d_{\max} = 4$ no es arbitraria: un Prisma
$\mathcal{P}(u,v,W)$ tiene di\'ametro de Hasse exactamente~2
($u \prec^{\!*} w \prec^{\!*} v$ para cualquier $w \in W$), y la
detecci\'on de amenazas~$K_5$ requiere explorar un salto adicional en
cada direcci\'on. As\'i, $d_{\max} = 2 + 2 = 4$ es la profundidad
m\'inima y suficiente para detectar toda la f\'isica local. Este valor
establece la \textbf{localidad fundamental} del C\'alculo de
Kuratowski: la f\'isica es local no por postulado, sino porque la
topolog\'ia del DAG solo requiere informaci\'on a 4 saltos de
profundidad.

\begin{result}[title={Constantes Topol\'ogicas, no Experimentales}]
\phantomsection\label{res:topological-constants}
Los tres par\'ametros $(\mathcal{V}_{\max},\; K_{\text{threat}},\;
d_{\max}) = (15,\; 2,\; 4)$ no son valores libres ajustados a datos:
son consecuencias de la estructura combinatoria del DAG causal
(estad\'istica de Poisson, planaridad de Kuratowski, y di\'ametro del
Prisma). Cualquier universo discreto, causal y planar que obedezca los
axiomas del C\'alculo de Kuratowski producir\'a necesariamente estos
mismos valores.
\end{result}

% ═══════════════════════════════════════════════════════════════════════════════
\chapter{La Integral Topol\'ogica: Acumulaci\'on y Volumen}
% ═══════════════════════════════════════════════════════════════════════════════
% CORRESPONDENCIA CENTRAL:
%   Continuo: ∫ d⁴x √(-g)  (integral de volumen sobre la variedad)
%   Discreto: |V ∩ ◇(A,B)|  (cardinalidad de eventos en un diamante)

\begin{introduction}
  \item El Diamante de Alexandrov como dominio natural de integraci\'on.
  \item La integral de volumen como cardinalidad:
        $\int d^4x\,\sqrt{-g} \;\to\; N\,\lP^4$.
  \item El ruido poissoniano de la integral y la Constante Cosmol\'ogica.
  \item Invarianza de Lorentz emergente del conteo discreto.
\end{introduction}

Si el diferencial de Kuratowski ($dx \to \prec^{\!*}$) nos dicta c\'omo
dar un paso indivisible en el espaciotiempo discreto, la operaci\'on
integral nos dicta c\'omo medir su inmensidad. En el c\'alculo
tradicional, la integral es el \'area bajo una curva suave o el volumen
de una variedad continua. En nuestra topolog\'ia fundamental, la
integraci\'on abandona la geometr\'ia continua y regresa a su forma
aritm\'etica m\'as pura: el conteo estricto.

% ─────────────────────────────────────────────────────────────────────────────
\section{El Diamante de Alexandrov y la Integral de Volumen}

Para integrar sobre una regi\'on del espaciotiempo, primero debemos
delimitarla utilizando \'unicamente relaciones causales, sin recurrir
a cintas m\'etricas externas.

\begin{definition}[Diamante de Alexandrov~\cite{alexandrov1967,kronheimer1967}]
\label{def:alexandrov}
Sean $A, B \in V$ dos eventos con $A \prec B$. El \textbf{Diamante de
Alexandrov} (o \textbf{intervalo causal}) es el conjunto de todos los
eventos causalmente intermedios:
\begin{equation}
  \Diamond(A,B) \;=\; \bigl\{\, e \in V : A \prec e \prec B \,\bigr\}.
  \label{eq:alexandrov}
\end{equation}
Por la finitud local del conjunto causal
(Definici\'on~\ref{def:event}), $\Diamond(A,B)$ es siempre un conjunto
finito. Su cardinalidad $|\Diamond(A,B)|$ est\'a bien definida.
\end{definition}

\noindent
El Diamante de Alexandrov es el dominio de integraci\'on can\'onico
del C\'alculo de Kuratowski. Toda regi\'on del espaciotiempo
delimitable por relaciones causales puede descomponerse en uni\'on de
diamantes. En la Relatividad General, el volumen espaciotemporal de
esta regi\'on se calcula mediante la integral:
\begin{equation}
  V_{\text{continuo}} \;=\;
  \int_{\Diamond} \sqrt{-g}\; d^4x.
  \label{eq:volume-continuum}
\end{equation}
En el C\'alculo de Kuratowski, esta integral se reemplaza por un
conteo:

\begin{definition}[Integral Topol\'ogica de Volumen]
\label{def:topological-integral}
Sea $\mathcal{C} = (V, \prec)$ un conjunto causal con densidad
fundamental de \emph{sprinkling}
$\rho = \lP^{-4}$ (un evento por volumen de Planck). La
\textbf{integral topol\'ogica de volumen} sobre un
Diamante de Alexandrov es:
\begin{equation}
  \boxed{\;
  \int_{\Diamond} \sqrt{-g}\; d^4x
  \quad\longrightarrow\quad
  \frac{|\Diamond(A,B)|}{\rho}
  \;=\; N\,\lP^4,
  \;}
  \label{eq:volume-discrete}
\end{equation}
donde $N = |\Diamond(A,B)|$ es el n\'umero de eventos encerrados en
el intervalo causal. La flecha denota convergencia estad\'istica en
el l\'imite $N \gg 1$ (v\'ease el Teorema~\ref{thm:volume-convergence}).
\end{definition}

\stoneref{$|\Diamond(A,B)|$ es el n\'umero de teoremas en la clausura
deductiva~\cite{tarski1930} entre las proposiciones $\varphi_A$
y~$\varphi_B$. Integrar sobre un volumen es contar los teoremas de
una teor\'ia.}

\noindent
Esta equivalencia es profunda: el volumen no es una propiedad
m\'etrica del espacio, sino una propiedad de \emph{conteo} de la
red de informaci\'on. Integrar es contar.

\begin{theorem}[Convergencia Volumen--Cardinalidad~\cite{myrheim1978,bombelli1987}]
\label{thm:volume-convergence}
Sea $\mathcal{C}$ un conjunto causal obtenido por \emph{sprinkling}
de Poisson con densidad~$\rho$ en un espaciotiempo globalmente
hiperb\'olico de dimensi\'on~$d$. Para cualesquiera $A \prec B$, la
variable aleatoria $N = |\Diamond(A,B)|$ satisface:
\begin{enumerate}
  \item \textbf{Valor esperado:}
    $\;\mathbb{E}[N] = \rho\,\mathrm{Vol}\bigl(\Diamond(A,B)\bigr).$
  \item \textbf{Fluctuaci\'on poissoniana:}
    $\;\mathrm{Var}[N] = \mathbb{E}[N],$
    de modo que $\sigma_N = \sqrt{N}$.
  \item \textbf{Convergencia relativa:}
    $\;\displaystyle\frac{\sigma_N}{\mathbb{E}[N]}
    = \frac{1}{\sqrt{N}}
    \;\xrightarrow{\;N \to \infty\;}\; 0.$
\end{enumerate}
\end{theorem}

\begin{proof}
Por construcci\'on, el \emph{sprinkling} de Poisson asigna a cada
regi\'on $R$ del espaciotiempo un n\'umero de eventos distribuido como
$\mathrm{Poisson}(\rho\,\mathrm{Vol}(R))$. El Diamante de Alexandrov
es una regi\'on medible del espaciotiempo, por lo que
$N \sim \mathrm{Poisson}(\rho\,\mathrm{Vol}(\Diamond))$.
Los puntos~(1) y~(2) son propiedades inmediatas de la distribuci\'on
de Poisson ($\mathrm{Var} = \mu$). El punto~(3) se sigue de la Ley
de los Grandes N\'umeros: para $N \gg 1$, la fracci\'on
$N/\mathbb{E}[N] \to 1$ casi seguramente, y el error relativo
$1/\sqrt{N}$ se desvanece.
\end{proof}

% ─────────────────────────────────────────────────────────────────────────────
\section{El Ruido de la Integral y la Constante Cosmol\'ogica ($\Lambda$)}

Redefinir la integral macrosc\'opica como un proceso de conteo
microsc\'opico tiene consecuencias f\'isicas profundas para los
misterios de la cosmolog\'ia moderna.

Como el conteo de~$N$ eventos en un conjunto causal aleatorio obedece
la estad\'istica de Poisson, la integral del volumen no es un valor
absolutamente est\'atico, sino que posee una incertidumbre
estad\'istica inherente. El volumen del universo tiene una fluctuaci\'on
fundamental:
\begin{equation}
  \delta V \;\sim\; \sqrt{N}\;\lP^4.
  \label{eq:volume-fluctuation}
\end{equation}

Esta ``vibraci\'on'' en el resultado de la integral no es un error de
medici\'on; es una propiedad ontol\'ogica del espaciotiempo discreto.
La pregunta es: ?`tiene consecuencias observables?

\begin{theorem}[Predicci\'on de Sorkin~\cite{sorkin1991,ahmed2004}]
\label{thm:sorkin-lambda}
Sea $\mathcal{C}$ un conjunto causal con $N$ eventos en el
4-volumen del universo observable, con densidad fundamental
$\rho = \lP^{-4}$. En el marco de la gravedad
unimodular~\cite{einstein1919,unruh1989}, donde la
constante cosmol\'ogica~$\Lambda$ es la variable conjugada al
4-volumen, la fluctuaci\'on poissoniana del conteo de eventos induce
una fluctuaci\'on efectiva en~$\Lambda$ de orden:
\begin{equation}
  \boxed{\;
  \Lambda_{\text{eff}}
  \;\sim\; \frac{1}{\sqrt{N}\;\lP^2}
  \;\sim\; \frac{\sqrt{\rho}}{\sqrt{N}}
  \;\sim\; H_0^2,
  \;}
  \label{eq:sorkin-lambda}
\end{equation}
donde $H_0$ es la constante de Hubble actual.
\end{theorem}

\begin{proof}
El volumen total del espaciotiempo observable es
$V_{\mathrm{obs}} = N/\rho = N\,\lP^4$, y su fluctuaci\'on
poissoniana es $\delta V = \sqrt{N}\,\lP^4$
(Teorema~\ref{thm:volume-convergence}). En gravedad unimodular, la
relaci\'on de incertidumbre volumen--constante cosmol\'ogica establece
$\delta\Lambda \cdot \delta V \sim 1$ (en unidades naturales
$\hbar = c = 1$). Por tanto:
\[
  \delta\Lambda
  \;\sim\; \frac{1}{\delta V}
  \;=\; \frac{1}{\sqrt{N}\,\lP^4}
  \;=\; \frac{\lP^{-4}}{\sqrt{N}}
  \;=\; \frac{\rho}{\sqrt{N}}.
\]
Ahora, $\Lambda$ tiene dimensiones de $[\text{longitud}]^{-2}$.
Reescribiendo con $\rho = \lP^{-4}$:
\[
  \delta\Lambda
  \;=\; \frac{1}{\sqrt{N}\,\lP^4}
  \;=\; \frac{1}{\lP^2} \cdot \frac{1}{\sqrt{N}\,\lP^2}
  \;=\; \frac{1}{\lP^2\,\sqrt{N\,\lP^4}}
  \;=\; \frac{1}{\lP^2 \sqrt{V_{\mathrm{obs}}}}.
\]
Para el universo observable, $V_{\mathrm{obs}} \sim (H_0^{-1})^4$
en unidades de Planck, de donde
$\sqrt{V_{\mathrm{obs}}} \sim H_0^{-2}/\lP^2$ y:
\[
  \delta\Lambda
  \;\sim\; \frac{1}{\lP^2} \cdot
           \frac{\lP^2}{H_0^{-2}}
  \;=\; H_0^2.
\]
Este valor coincide, en orden de magnitud, con el valor observado
$\Lambda_{\text{obs}} \approx 10^{-122}\,\lP^{-2}$.
\end{proof}

\stoneref{$\Lambda$ es la fluctuaci\'on de Poisson en el n\'umero de
teoremas del sistema deductivo cosmol\'ogico: el \emph{ruido de fondo
l\'ogico} de la clausura deductiva~\cite{tarski1930}. La constante
cosmol\'ogica es el residuo estad\'istico de contar teoremas.}

\noindent
La energ\'ia que acelera la expansi\'on del cosmos no es un fluido
ex\'otico inyectado en las ecuaciones de campo; es, anal\'iticamente,
el residuo estad\'istico de la Integral de Kuratowski al operar sobre
un n\'umero inmenso, pero finito, de eventos causales.

\begin{result}[title={La Constante Cosmol\'ogica como Ruido de Conteo}]
\phantomsection\label{res:cosmological-constant}
La predicci\'on de Sorkin precede en siete a\~nos al descubrimiento
experimental de la expansi\'on
acelerada~\cite{perlmutter1999,riess1998}.
El C\'alculo de Kuratowski la hace inevitable: si la integral de
volumen es un conteo de Poisson, la constante cosmol\'ogica \emph{no
puede ser cero ni infinita}. Su valor est\'a fijado por la
cardinalidad del universo observable:
\[
  \Lambda \;\sim\; N^{-1/2}\,\lP^{-2},
  \qquad
  N \;\sim\; 10^{240}.
\]
El problema de la constante cosmol\'ogica~\cite{weinberg1989} se
disuelve: $\Lambda$ no es ni cero ni $\lP^{-2}$, sino el residuo
estad\'istico del conteo
causal~\cite{sorkin2005lambda,zwane2018}.
\end{result}

% ─────────────────────────────────────────────────────────────────────────────
\section{Invarianza de Lorentz desde el Conteo}

La invarianza de Lorentz del C\'alculo de Kuratowski no se impone
como postulado: emerge del conteo. El argumento descansa sobre un
\'unico hecho algebraico y una consecuencia probabil\'istica.

\begin{theorem}[Bombelli--Henson--Sorkin~\cite{bombelli2009}]
\label{thm:lorentz-sprinkling}
Sea $\Lambda$ una transformaci\'on de Lorentz y sea $\mathcal{C}$ un
conjunto causal obtenido por \emph{sprinkling} de Poisson con
densidad~$\rho$ en un espaciotiempo lorentziano. Entonces la
distribuci\'on de~$\mathcal{C}$ es estad\'isticamente invariante
bajo~$\Lambda$: el n\'umero de eventos en cualquier
Diamante~$\Diamond(A,B)$ es un escalar de Lorentz.
\end{theorem}

\begin{proof}
Una transformaci\'on de Lorentz satisface $\det\Lambda = 1$ por
pertenencia a $SO^+(1,3)$. Por lo tanto, preserva el elemento de
4-volumen lorentziano $d^4x\,\sqrt{-g}$. El n\'umero de eventos
esparcidos en una regi\'on~$R$ se distribuye como
$\mathrm{Poisson}(\rho\,\mathrm{Vol}(R))$
(Teorema~\ref{thm:volume-convergence}). Puesto que
$\mathrm{Vol}(\Lambda R) = \mathrm{Vol}(R)$, la variable aleatoria
$N = |\Diamond(A,B) \cap V|$ tiene id\'entica distribuci\'on en todo
marco inercial. En particular, su valor esperado, su varianza y
todos sus momentos son escalares de Lorentz.
\end{proof}

\noindent
La consecuencia para la integral topol\'ogica es directa: dado que
la cardinalidad~$N$ es invariante, el volumen discreto
$N\,\lP^4$ (Definici\'on~\ref{def:topological-integral}) tambi\'en
lo es. La invarianza de Lorentz que el c\'alculo continuo postula
para la integral $\int\!\sqrt{-g}\,d^4x$ no es un axioma
independiente; es una propiedad heredada del conteo de eventos en un
\emph{sprinkling} de Poisson. El \emph{sprinkling} es, de hecho, la
\'unica distribuci\'on puntual sobre un espaciotiempo lorentziano que
respeta esta invarianza~\cite{bombelli2009}.

% ═══════════════════════════════════════════════════════════════════════════════
\chapter{La Derivada Causal: Difusi\'on y Resoluci\'on}
% ═══════════════════════════════════════════════════════════════════════════════
% CORRESPONDENCIA CENTRAL:
%   Continuo: ∂f/∂x^μ, ∇_μ, □  (derivadas, conexión, d'Alembertiano)
%   Discreto: difusión sobre el Hasse, valencia local, flujo neto

\begin{introduction}
  \item El operador de difusi\'on y la dimensi\'on espectral como derivada.
  \item El Teorema de la Resoluci\'on Causal: precisi\'on de la derivada
        discreta.
  \item Gradientes de flujo, inercia y ecuaciones de continuidad.
\end{introduction}

Si la integral de Kuratowski nos da el volumen (el conteo de eventos),
la derivada nos revela la dimensionalidad y la curvatura. En este
cap\'itulo, definimos el operador diferencial no como una resta de
valores en puntos adyacentes, sino como la tasa de propagaci\'on de
la informaci\'on a trav\'es de la red causal.

% ─────────────────────────────────────────────────────────────────────────────
\section{El Operador de Difusi\'on y la Dimensi\'on Espectral}

En el c\'alculo de Riemann, la dimensi\'on de una variedad est\'a dada
de forma aprior\'istica: se fija $d = 4$ antes de escribir una sola
ecuaci\'on. En el C\'alculo de Kuratowski, la dimensi\'on es una
propiedad
\emph{emergente}~\cite{carlip2017,ambjorn2005,eichhorn2014},
detectable mediante un operador de difusi\'on (el caminante aleatorio
sobre el diagrama de Hasse).

\begin{definition}[Caminante Aleatorio sobre el Diagrama de Hasse]
\label{def:random-walker}
Sea $H = (V, E)$ el diagrama de Hasse (no dirigido) de un conjunto
causal $\mathcal{C} = (V, \prec)$, donde $E$ es el conjunto de
v\'inculos de cobertura~$\prec^{\!*}$ con la orientaci\'on suprimida.
Un \textbf{caminante aleatorio} parte de un evento~$u_0 \in V$ y, en
cada paso discreto $t \to t+1$, salta a un vecino elegido
uniformemente al azar:
\[
  \Pr\bigl[u_{t+1} = v \;\big|\; u_t = u\bigr]
  \;=\;
  \begin{cases}
    1/\deg(u) & \text{si } (u,v) \in E,\\[2pt]
    0 & \text{en caso contrario,}
  \end{cases}
\]
donde $\deg(u) = |\{v : (u,v) \in E\}|$ es la valencia total
(in-degree $+$ out-degree en el DAG original).
\end{definition}

\begin{definition}[Probabilidad de Retorno y Dimensi\'on Espectral]
\label{def:spectral-dimension}
La \textbf{probabilidad de retorno} $P(t)$ es la probabilidad de que
el caminante regrese a su origen despu\'es de $t$~pasos, promediada
sobre todos los or\'igenes posibles:
\[
  P(t) \;=\; \frac{1}{|V|}\sum_{u \in V}
  \Pr\bigl[u_t = u \;\big|\; u_0 = u\bigr].
\]
La \textbf{dimensi\'on espectral} se extrae como la pendiente
logar\'itmica de esta probabilidad:
\begin{equation}
  \boxed{\;
  \dS(t) \;=\; -2\,\frac{d\log P(t)}{d\log t}.
  \;}
  \label{eq:spectral-dim}
\end{equation}
Para una variedad suave $d$-dimensional, $P(t) \sim t^{-d/2}$
y por tanto $\dS(t) \to d$.
\end{definition}

\stoneref{$\dS$ es la \emph{dimensionalidad inferencial}: el n\'umero
de ejes de inferencia independientes del sistema deductivo.
$\dS \to 4$ significa que existen exactamente cuatro direcciones
l\'ogicas ortogonales en las que el sistema puede razonar. La
pregunta ``?`por qu\'e 4 dimensiones?'' se convierte en ``?`por qu\'e
4 ejes de inferencia independientes?''}

\noindent
La dimensi\'on espectral es la \emph{derivada} fundamental del
C\'alculo de Kuratowski: mide c\'omo decae la propagaci\'on de la
informaci\'on con la profundidad temporal. Es el an\'alogo discreto
de la traza del n\'ucleo del calor
$K(t) = \mathrm{Tr}\,e^{-t\Delta}$ en geometr\'ia espectral
riemanniana~\cite{minakshisundaram1949,gilkey1975}. La
dimensi\'on espectral fue introducida como diagn\'ostico de
conjuntos causales en~\cite{eichhorn2014}; el fen\'omeno de
reducci\'on dimensional a escalas cortas ($\dS \to 2$) aparece
tambi\'en en triangulaciones din\'amicas
causales~\cite{ambjorn2005} y otros enfoques de gravedad
cu\'antica~\cite{carlip2009,modesto2009}.

\begin{result}[title={Estabilizaci\'on Dimensional por la Materia}]
\phantomsection\label{res:dimensional-stabilization}
Los resultados num\'ericos de la simulaci\'on
(\texttt{prism\_simulation}, $N$ hasta $5 \times 10^6$) revelan un
fen\'omeno consistente: mientras el vac\'io puro exhibe una dimensi\'on
espectral ruidosa y fractal, la presencia de Prismas
Causales~$K_{2,N}$ \textbf{estabiliza} $\dS$ en torno a~$4$ a tiempos
de difusi\'on intermedios. La materia no es un hu\'esped del
espaciotiempo tetradimensional; es, por definici\'on operacional, la
estructura topol\'ogica que \emph{fuerza} a la derivada de la
informaci\'on a comportarse tetradimensionalmente.
\end{result}

% ─────────────────────────────────────────────────────────────────────────────
\section{El Teorema de la Resoluci\'on Causal}

La validez de cualquier derivada discreta depende de la resoluci\'on
del sondeo. En el C\'alculo de Kuratowski, la precisi\'on de nuestra
descripci\'on del espaciotiempo est\'a intr\'insecamente ligada al
n\'umero de caminos causales sondeados.

El Teorema~\ref{thm:causal-resolution} (Cap\'itulo~6) establece que
para mantener una tolerancia de error~$\epsilon$ en la medici\'on de
la geometr\'ia hasta una profundidad~$t_{\max}$, el n\'umero de
caminantes debe satisfacer:
\begin{equation}
  W \;\geq\; \frac{t_{\max}^{\,2}}{c\;\epsilon^2}.
  \tag{\ref{eq:causal-resolution}}
\end{equation}

\noindent
Este resultado establece el \textbf{l\'imite de certidumbre}:
la ilusi\'on del continuo de Riemann solo es sostenible si la densidad
de interacciones causales es lo suficientemente alta como para
promediar el ruido de Poisson de la topolog\'ia discreta. El escalado
$W \propto t_{\max}^2$ es una consecuencia directa de $\dS = 4$:
en un espaciotiempo de dimensi\'on espectral~$d$, la cota general
es $W \propto t_{\max}^{d/2}$.

\begin{theorem}[Cota General de Resoluci\'on Espectral]
\label{thm:resolution-general}
Para un conjunto causal con dimensi\'on espectral~$\dS$ a tiempos
grandes, $P(t) \sim t^{-\dS/2}$. La cota de resoluci\'on del
Teorema~\ref{thm:causal-resolution} se generaliza a:
\begin{equation}
  W \;\geq\; \frac{t_{\max}^{\,\dS/2}}{c\;\epsilon^2}.
  \label{eq:resolution-general}
\end{equation}
En particular:
\begin{enumerate}
  \item $\dS = 2$ (l\'imite UV, grafos tipo \'arbol):
        $W \propto t_{\max}$\quad (costo lineal).
  \item $\dS = 4$ (espaciotiempo macrosc\'opico):
        $W \propto t_{\max}^2$\quad (costo cuadr\'atico).
  \item $\dS > 4$ (dimensiones extra compactas):
        el costo crece m\'as r\'apido que cuadr\'atico.
\end{enumerate}
\end{theorem}

\begin{proof}
Sigue id\'enticamente la prueba del
Teorema~\ref{thm:causal-resolution}, sustituyendo
$P(t_{\max}) = c\,t_{\max}^{-\dS/2}$:
\[
  W \;\geq\; \frac{1}{\epsilon^2\,P(t_{\max})}
  \;=\; \frac{t_{\max}^{\,\dS/2}}{c\;\epsilon^2}.
\]
\end{proof}

\noindent
La f\'ormula~\eqref{eq:resolution-general} revela que la
dimensi\'on espectral controla directamente el
\emph{precio computacional de la suavidad}: dimensiones m\'as altas
son m\'as costosas de sostener.

% ─────────────────────────────────────────────────────────────────────────────
\section{Anticadenas: Los Cortes Espaciales del DAG}

Para definir gradientes espaciales sobre la red causal, necesitamos
la noci\'on de ``instante discreto'': un corte del DAG que separe
pasado de futuro.

\begin{definition}[Anticadena]
\label{def:antichain}
Una \textbf{anticadena} es un subconjunto $\mathcal{A} \subset V$
cuyos elementos son mutuamente incomparables bajo el orden causal:
\[
  \forall\, a, b \in \mathcal{A},\;a \neq b
  \;\colon\quad
  a \nprec b \;\;\text{y}\;\; b \nprec a.
\]
Una anticadena es el an\'alogo discreto de una
\textbf{hipersuperficie de simultaneidad}: define un ``instante''
donde ning\'un evento puede influenciar causalmente a otro.
\end{definition}

\noindent
El DAG causal admite una \textbf{foliaci\'on en estratos de
Hasse}~\cite{rideout2000,surya2019}: una partici\'on
$V = \mathcal{A}_0 \cup \mathcal{A}_1 \cup \cdots \cup
\mathcal{A}_T$ donde cada estrato~$\mathcal{A}_k$ es una anticadena
maximal, y todo v\'inculo $u \prec^{\!*} v$ conecta un
estrato~$\mathcal{A}_k$ con el estrato siguiente~$\mathcal{A}_{k+1}$.
Los estratos definen un ``tiempo discreto global'' $k = 0, 1,
\ldots, T$.

% ─────────────────────────────────────────────────────────────────────────────
\section{Gradientes de Flujo e Inercia}

La derivada de Kuratowski nos permite definir la fuerza no como un
campo vectorial externo, sino como un \emph{gradiente de
conectividad}.

\begin{definition}[Operador d'Alembertiano
Discreto~\cite{sorkin2007box,benincasa2010,dowker2013}]
\label{def:box-operator}
Sea $\phi \colon V \to \mathbb{R}$ una funci\'on escalar sobre los
eventos del conjunto causal. El \textbf{d'Alembertiano discreto}
(u operador de onda causal) en un evento~$u$ se define como la
diferencia entre el flujo saliente y el flujo entrante a trav\'es de
los v\'inculos de cobertura:
\begin{equation}
  (\Box_{\mathcal{C}}\,\phi)(u)
  \;=\;
  \underbrace{\sum_{v :\, u \prec^* v} \phi(v)}_{\text{hijos de } u}
  \;-\;
  \underbrace{\sum_{w :\, w \prec^* u} \phi(w)}_{\text{padres de } u}.
  \label{eq:box-discrete}
\end{equation}
\end{definition}

\noindent
Este operador tiene una interpretaci\'on transparente. Sean
$d^+(u) = |\{v : u \prec^{\!*} v\}|$ el out-degree y
$d^-(u) = |\{w : w \prec^{\!*} u\}|$ el in-degree de~$u$. Entonces:
\begin{itemize}
  \item Si $\phi \equiv 1$ (campo constante):
    $(\Box_{\mathcal{C}}\,\phi)(u) = d^+(u) - d^-(u)$.
    El operador mide la \textbf{asimetr\'ia causal local}---la
    diferencia entre el n\'umero de futuros inmediatos y pasados
    inmediatos.
  \item Si $\phi$ var\'ia, el operador detecta c\'omo se
    \textbf{propaga} la se\~nal: la suma sobre hijos mide la
    dispersi\'on hacia el futuro, y la suma sobre padres mide la
    convergencia desde el pasado.
\end{itemize}

\noindent
A partir de~\eqref{eq:box-discrete}, el \textbf{gradiente causal
normalizado} en la direcci\'on temporal es:

\begin{definition}[Gradiente Causal]
\label{def:causal-gradient}
\begin{equation}
  (\nabla_t\,\phi)(u)
  \;=\;
  \frac{1}{d^+(u)}\sum_{v :\, u \prec^* v} \phi(v)
  \;-\;
  \frac{1}{d^-(u)}\sum_{w :\, w \prec^* u} \phi(w).
  \label{eq:causal-gradient}
\end{equation}
Es el promedio hacia el futuro menos el promedio desde el pasado: la
tasa de cambio discreta de~$\phi$ por salto topol\'ogico.
\end{definition}

\subsection{Inercia como Resistencia al Flujo}

Cuando un caminante atraviesa un Prisma Causal~$K_{2,N}$, experimenta
una disminuci\'on en su velocidad de propagaci\'on efectiva debido a la
bifurcaci\'on de caminos en los $N$ nodos intermedios. Formalicemos
esta observaci\'on.

\begin{theorem}[Retardo del Gradiente por Prismas Causales]
\label{thm:gradient-delay}
Sea $\phi$ un campo escalar que se propaga a trav\'es de un
Prisma~$\mathcal{P}(u,v,W)$ con $|W| = N$ nodos intermedios
(Definici\'on~\ref{def:causal-prism}, Vol.~II). Para el polo de
entrada~$u$, el gradiente causal satisface:
\begin{equation}
  (\nabla_t\,\phi)(u)
  \;=\;
  \frac{1}{N}\sum_{w \in W} \phi(w)
  \;-\;
  \frac{1}{d^-(u)}\sum_{w' :\, w' \prec^* u} \phi(w'),
  \label{eq:gradient-prism}
\end{equation}
donde la primera suma se distribuye uniformemente sobre los $N$
intermedios. Si $\phi$ es constante sobre~$W$ (campo homog\'eneo en
el vientre del prisma), entonces:
\[
  (\nabla_t\,\phi)(u)
  \;=\; \phi_W - \bar{\phi}_{\text{pasado}},
\]
independientemente de~$N$. Sin embargo, la informaci\'on requiere
$N$~pasos de difusi\'on para atravesar el vientre antes de alcanzar el
polo de salida~$v$, produciendo un \textbf{retardo efectivo}
$\Delta t_{\text{prisma}} = N$.
\end{theorem}

\begin{proof}
En el Prisma $\mathcal{P}(u,v,W)$, el polo~$u$ tiene $d^+(u) = N$
hijos (todos los $w \in W$). Por la
Definici\'on~\ref{def:causal-gradient}:
\[
  (\nabla_t\,\phi)(u)
  \;=\;
  \frac{1}{N}\sum_{w \in W}\phi(w) - \bar{\phi}_{\text{pasado}}.
\]
Un caminante que entra en~$u$ debe elegir uno de los $N$ caminos
intermedios. Cada intermedio~$w_i$ est\'a conectado a~$v$ por un
v\'inculo de cobertura $w_i \prec^{\!*} v$. Pero antes de que la
informaci\'on pueda propagarse de~$u$ a~$v$, debe pasar por al menos
un intermedio: el camino m\'inimo es
$u \prec^{\!*} w_i \prec^{\!*} v$, de longitud~2. Comparado con un
v\'inculo de vac\'io directo $u \prec^{\!*} v'$ (longitud~1), el
Prisma introduce un retardo de $+1$ paso por cada intermedio
visitado. Si el caminante explora los $N$ intermedios por difusi\'on,
el tiempo de primer paso promedio al polo~$v$ escala como $\Delta
t_{\text{prisma}} \sim N$.
\end{proof}

\noindent
Esta ``resistencia'' al flujo de informaci\'on es la representaci\'on
matem\'atica de la inercia. El objeto no ``pesa'' porque posea una
propiedad intr\'inseca misteriosa; pesa porque la derivada de su
avance causal es retardada por la topolog\'ia del Prisma.

\begin{result}[title={La Derivada de la Inercia}]
\phantomsection\label{res:inertia-derivative}
La masa topol\'ogica $M = \kappa N$
(Axioma de Inercia Topol\'ogica, Vol.~II) se manifiesta
din\'amicamente como el retardo del gradiente causal:
\[
  \Delta t_{\text{prisma}} \;\sim\; N \;=\; M/\kappa.
\]
A mayor masa, mayor retardo. La inercia es, literalmente, el n\'umero
de v\'inculos de cobertura que la informaci\'on debe procesar antes de
avanzar un paso global.
\end{result}

% ─────────────────────────────────────────────────────────────────────────────
% TODO — §3.5: Ecuaciones de Continuidad en el Grafo
% Conservación de flujo: lo que entra por los padres sale por los hijos.
% Definición del operador divergencia discreto.
% La ecuación de continuidad discreta: Σ_{in} J - Σ_{out} J = 0.
% Límite continuo: ∂_μ J^μ = 0.
% Conexión con los random walkers de la dimensión espectral (prism_simulation).
% Refs: \cite{sorkin2012,johnston2009,aslanbeigi2014}

% ███████████████████████████████████████████████████████████████████████████████
%   PART II : LA FÍSICA EMERGENTE
% ███████████████████████████████████████████████████████████████████████████████
\part{La F\'isica Emergente}

% ═══════════════════════════════════════════════════════════════════════════════
\chapter{Geod\'esicas y el Principio de M\'aximo Tiempo}
% ═══════════════════════════════════════════════════════════════════════════════
% CORRESPONDENCIA CENTRAL:
%   Continuo: δ∫L dτ = 0  (extremo del funcional de acción)
%   Discreto: cadena causal más larga entre A y B (tiempo propio máximo)

\begin{introduction}
  \item La geod\'esica como cadena causal de longitud m\'axima.
  \item Inercia y desviaci\'on geod\'esica: el Prisma como obst\'aculo
        topol\'ogico.
  \item El principio de m\'inima acci\'on l\'ogica: eficiencia causal.
\end{introduction}

En la geometr\'ia continua de Riemann~\cite{wald1984}, una
part\'icula libre sigue una geod\'esica, definida como la curva que
extremiza la distancia entre dos puntos. En signatura euclidiana, el
extremo es un m\'inimo. Sin embargo, en la signatura lorentziana del
espaciotiempo, el principio se invierte: la geod\'esica temporal
\textbf{maximiza} el tiempo propio~\cite{mtw1973}. El C\'alculo de Kuratowski formaliza esta inversi\'on de
manera natural mediante la estructura de cadenas del DAG.

% ─────────────────────────────────────────────────────────────────────────────
\section{La Geod\'esica como Cadena M\'axima}

En el Cap\'itulo~1, definimos el tiempo propio a lo largo de una
cadena $\gamma$ como su n\'umero de v\'inculos de cobertura:
$\tau(\gamma) = |\gamma|$
(Definici\'on~\ref{def:proper-time}). El Teorema de Myrheim--Meyer
(\ref{thm:myrheim-meyer}) demostr\'o que, en el l\'imite de alta
densidad, la cadena m\'as larga converge al tiempo propio
geod\'esico del continuo. Ahora formalizamos la geod\'esica discreta
como objeto central:

\begin{definition}[Geod\'esica de Kuratowski]
\label{def:geodesic}
Sean $A, B \in V$ con $A \prec B$. La \textbf{geod\'esica de
Kuratowski} entre~$A$ y~$B$ es cualquier cadena causal que alcanza
la longitud m\'axima. La \textbf{longitud geod\'esica} es:
\begin{equation}
  \boxed{\;
  \tau(A, B)
  \;=\;
  \max_{\gamma \,\in\, \Gamma(A,B)} |\gamma|
  \;=\; \ell(A,B),
  \;}
  \label{eq:geodesic-length}
\end{equation}
donde $\Gamma(A,B)$ es el conjunto de todas las cadenas entre~$A$
y~$B$ (Definici\'on~\ref{def:chain}).
\end{definition}

\stoneref{La geod\'esica es la \emph{demostraci\'on \'optima}: la
cadena inferencial que extrae el m\'aximo contenido deductivo entre
las premisas $\varphi_A$ y~$\varphi_B$. La naturaleza elige la
prueba m\'as detallada.}

\noindent
Obs\'ervese la inversi\'on respecto a la intuici\'on euclidiana: la
geod\'esica no minimiza, sino que \textbf{maximiza}. En el espacio
euclidiano, el camino m\'as corto gana. En el espaciotiempo causal,
el camino m\'as largo gana---porque acumula m\'as v\'inculos de
cobertura y, por tanto, m\'as tiempo propio. Esto es exactamente el
principio de Relatividad General: entre dos eventos, el observador
inercial (en ca\'ida libre) es aquel cuyo reloj marca el
\emph{mayor} intervalo de tiempo.

\begin{theorem}[Unicidad Asint\'otica de la
Geod\'esica~\cite{brightwell1991}]
\label{thm:geodesic-uniqueness}
Sea $\mathcal{C}$ un conjunto causal obtenido por \emph{sprinkling}
de Poisson con densidad~$\rho$ en un espaciotiempo globalmente
hiperb\'olico. Para $A \prec B$ con $\ell(A,B) \gg 1$, la cadena
m\'axima es asint\'oticamente \'unica en el siguiente sentido: si
$\gamma_1$ y $\gamma_2$ son dos cadenas de longitud~$\ell(A,B)$,
entonces sus elementos difieren en a lo sumo $O(\sqrt{\ell})$ nodos.
\end{theorem}

\begin{proof}
Sea $\gamma^* = (e_0, e_1, \ldots, e_\ell)$ una cadena m\'axima.
Cualquier cadena alternativa~$\gamma'$ de igual longitud debe pasar
por los mismos ``cuellos de botella'' del DAG---regiones donde el
diamante $\Diamond(A,B)$ se estrecha. En un \emph{sprinkling} de
Poisson $d$-dimensional, el n\'umero de eventos en una secci\'on
transversal a profundidad~$k$ fluctu\'a como $\sqrt{\rho\,\sigma_k}$,
donde $\sigma_k$ es el $(d{-}1)$-volumen de la secci\'on. Dos cadenas
m\'aximas solo pueden divergir donde la secci\'on transversal admite
m\'ultiples caminos de igual longitud, lo cual ocurre en una
fracci\'on $O(1/\sqrt{\ell})$ de los pasos. As\'i, el n\'umero de
nodos en que difieren es $O(\sqrt{\ell})$, y la fracci\'on de
discrepancia tiende a cero: $O(\sqrt{\ell})/\ell = O(1/\sqrt{\ell})
\to 0$.
\end{proof}

\noindent
Esta cuasi-unicidad explica por qu\'e la materia macrosc\'opica
experimenta el tiempo como algo flu\'ido y determinista: el promedio
sobre trillones de cadenas m\'aximas converge a una \'unica
geod\'esica continua, con fluctuaciones relativas imperceptibles.

% ─────────────────────────────────────────────────────────────────────────────
\section{Inercia y Desviaci\'on Geod\'esica}

?`Por qu\'e una part\'icula se siente ``pesada''? En el C\'alculo de
Kuratowski, la inercia es la resistencia de la topolog\'ia a permitir
cadenas largas.

\begin{definition}[Resistencia Topol\'ogica]
\label{def:topological-resistance}
Sea $\Diamond(A,B)$ un diamante que contiene una densidad
$\nu$ de Prismas Causales por unidad de profundidad causal. La
\textbf{resistencia topol\'ogica} del diamante es:
\begin{equation}
  \mathcal{R}(A,B)
  \;=\;
  \frac{\ell_{\text{vac}}(A,B) - \ell(A,B)}{\ell_{\text{vac}}(A,B)},
  \label{eq:topological-resistance}
\end{equation}
donde $\ell_{\text{vac}}(A,B)$ es la longitud geod\'esica que
tendr\'ia el mismo diamante en ausencia de Prismas (vac\'io puro), y
$\ell(A,B)$ es la longitud geod\'esica real.
\end{definition}

\noindent
El Teorema~\ref{thm:time-dilation} (Cap\'itulo~1) ya demostr\'o que
los Prismas consumen v\'inculos en rutas internas, reduciendo el
avance causal global. Ahora podemos cuantificar el efecto sobre la
geod\'esica:

\begin{theorem}[Acortamiento Geod\'esico por Prismas]
\label{thm:geodesic-shortening}
Sea $\Diamond(A,B)$ un diamante que contiene $p$~Prismas
$\mathcal{P}_1, \ldots, \mathcal{P}_p$ con
$N_j = |W_j|$ intermedios cada uno. La longitud geod\'esica efectiva
satisface:
\begin{equation}
  \ell(A,B)
  \;\leq\;
  \ell_{\text{vac}}(A,B) - \sum_{j=1}^{p} (N_j - 1),
  \label{eq:geodesic-shortening}
\end{equation}
donde el t\'ermino $N_j - 1$ refleja que cada Prisma desvía la
cadena m\'axima por al menos $N_j - 1$ v\'inculos adicionales que no
contribuyen al avance global.
\end{theorem}

\begin{proof}
En el vac\'io, la cadena m\'axima usa cada v\'inculo de cobertura para
avanzar causalmente de $A$ hacia~$B$. Cuando la cadena atraviesa el
Prisma $\mathcal{P}_j(u_j, v_j, W_j)$, el camino
$u_j \prec^{\!*} w \prec^{\!*} v_j$ consume 2~v\'inculos para
recorrer una distancia causal que, en el vac\'io, requerir\'ia 1
(el v\'inculo directo $u_j \prec^{\!*} v_j$ no existe porque $W_j$
lo intercepta). El costo neto por Prisma es al menos~$+1$ v\'inculo
por intermedio visitado. Si la cadena m\'axima debe atravesar los
$N_j$~intermedios secuencialmente (caso peor), el costo es $N_j - 1$
v\'inculos consumidos sin avance global. Sumando sobre todos los
Prismas:
\[
  \ell(A,B)
  \;\leq\;
  \ell_{\text{vac}}(A,B)
  \;-\; \sum_{j=1}^{p} (N_j - 1).
\]
\end{proof}

\noindent
La desviaci\'on geod\'esica---la divergencia entre dos cadenas
m\'aximas inicialmente cercanas---es el mecanismo por el cual la
curvatura se manifiesta. En una regi\'on de vac\'io uniforme
($\nu = 0$), dos geod\'esicas paralelas permanecen paralelas. En
una regi\'on con Prismas ($\nu > 0$), una geod\'esica que pasa cerca
de un Prisma se desv\'ia (pierde v\'inculos en procesamiento interno),
mientras que una geod\'esica distante no. Esta diferencia de longitud
es la \textbf{curvatura discreta}, an\'aloga a la ecuaci\'on de
desviaci\'on geod\'esica de Jacobi en el
continuo~\cite{wald1984,hawking1973}.

% ─────────────────────────────────────────────────────────────────────────────
\section{El Principio de M\'inima Acci\'on L\'ogica}

En la f\'isica cl\'asica, la acci\'on $S[\gamma] = \int_\gamma L\,d\tau$
es un funcional escalar sobre trayectorias, y las ecuaciones del
movimiento emergen de la condici\'on $\delta S = 0$. En el C\'alculo
de Kuratowski, reinterpretamos la acci\'on como una medida de
\textbf{eficiencia l\'ogica} del
grafo~\cite{benincasa2010,benincasa2011,buck2015}.

\begin{definition}[Funcional de Acci\'on Discreta]
\label{def:discrete-action}
Sea $\gamma \in \Gamma(A,B)$ una cadena causal. El \textbf{funcional
de acci\'on de Kuratowski} es:
\begin{equation}
  S[\gamma] \;=\; -|\gamma|.
  \label{eq:discrete-action}
\end{equation}
El signo negativo garantiza que \emph{minimizar}~$S$ equivale a
\emph{maximizar} la longitud de la cadena---recuperando el principio
de m\'aximo tiempo propio.
\end{definition}

\stoneref{$S[\gamma] = -|\gamma|$ es el \emph{principio de m\'axima
longitud de prueba}: la naturaleza elige la demostraci\'on m\'as
detallada. Minimizar la acci\'on equivale a maximizar el n\'umero de
pasos inferenciales at\'omicos---la prueba que no omite ning\'un
lema.}

\begin{theorem}[Condici\'on Extremal Discreta]
\label{thm:extremal-condition}
Una cadena $\gamma^* \in \Gamma(A,B)$ es un extremo del funcional de
acci\'on $S[\gamma] = -|\gamma|$ si y solo si no admite
\textbf{inserciones}: no existe ning\'un evento $e \in \Diamond(A,B)
\setminus \gamma^*$ tal que, al intercalarlo en~$\gamma^*$, la cadena
resultante sea m\'as larga.

Formalmente, $\gamma^*$ es extremal si y solo si $\gamma^*$ es una
\textbf{cadena m\'axima}:
\begin{equation}
  |\gamma^*| \;=\; \ell(A,B)
  \;=\; \max_{\gamma \in \Gamma(A,B)} |\gamma|.
  \label{eq:extremal}
\end{equation}
\end{theorem}

\begin{proof}
La ``variaci\'on'' discreta $\delta\gamma$ consiste en insertar un
evento~$e$ entre dos eslabones consecutivos $e_k$ y $e_{k+1}$ de la
cadena. Esto es posible si y solo si existen v\'inculos
$e_k \prec^{\!*} e \prec^{\!*} e_{k+1}$, en cuyo caso la longitud
aumenta en~$+1$. Si $\gamma^*$ admite tal inserci\'on, entonces
$|\gamma^*| < |\gamma^* \cup \{e\}|$, contradiciendo la maximalidad.
Rec\'iprocamente, si $\gamma^*$ es m\'axima, ninguna inserci\'on puede
aumentar su longitud.
\end{proof}

\noindent
Una part\'icula libre no ``elige'' su camino; el camino emerge de la
configuraci\'on del grafo que permite la mayor propagaci\'on de
causalidad. La geod\'esica de Kuratowski es el camino de menor
resistencia informativa, que resulta ser, parad\'ojicamente, el que
acumula la mayor cantidad de eventos de tiempo propio: el camino donde
la informaci\'on fluye con la menor cantidad de obstrucciones
topol\'ogicas ($K_5$).

% ═══════════════════════════════════════════════════════════════════════════════
\chapter{El Tensor de Inercia y la M\'etrica Emergente}
% ═══════════════════════════════════════════════════════════════════════════════
% CORRESPONDENCIA CENTRAL:
%   Continuo: T_μν (tensor energía-momento) y g_μν (métrica)
%   Discreto: Prisma Causal K_{2,N} y densidad de vínculos intermedios

\begin{introduction}
  \item El Prisma Causal como operador de retraso: el tensor de inercia local.
  \item La m\'etrica $g_{\mu\nu}$ emergente de la probabilidad de retorno.
  \item Invarianza de Lorentz en el discreto: de Poisson a la valencia.
  \item Gravedad como estad\'istica de caminos.
\end{introduction}

En la Relatividad General, la materia le dice al espacio c\'omo
curvarse a trav\'es del tensor de energ\'ia-impulso $T_{\mu\nu}$. En
el C\'alculo de Kuratowski, este proceso no es una interacci\'on entre
dos entidades distintas, sino una manifestaci\'on de la densidad de
informaci\'on. La curvatura es la respuesta estad\'istica del grafo a
la acumulaci\'on de defectos topol\'ogicos.

% ─────────────────────────────────────────────────────────────────────────────
\section{El Prisma Causal como Operador de Retraso}

Un Prisma Causal $K_{2,N}$ act\'ua como un \emph{colador informativo}.
Cuando una geod\'esica (cadena m\'axima) lo atraviesa, el n\'umero de
caminos posibles se expande, pero la conectividad local se vuelve
``pesada'': la cadena debe invertir v\'inculos en procesamiento
interno (Teorema~\ref{thm:geodesic-shortening}).

Para cuantificar este efecto, definimos la carga de inercia que cada
evento soporta:

\begin{definition}[Tensor de Inercia Local]
\label{def:inertia-tensor}
Sea $\mathcal{C} = (V, \prec)$ un conjunto causal que contiene un
conjunto de Prismas
$\{\mathcal{P}_k(u_k, v_k, W_k)\}_{k=1}^{p}$. El \textbf{tensor de
inercia local} en un evento~$u \in V$ es:
\begin{equation}
  \mathcal{I}(u)
  \;=\;
  \sum_{k=1}^{p} \mathbf{1}_{u \in \mathcal{P}_k} \;\cdot\; N_k,
  \label{eq:inertia-tensor}
\end{equation}
donde $\mathbf{1}_{u \in \mathcal{P}_k}$ es la funci\'on indicadora
($= 1$ si $u$ pertenece al Prisma~$\mathcal{P}_k$ como polo o
intermedio, $= 0$ en caso contrario) y $N_k = |W_k|$ es la masa
topol\'ogica del Prisma~$k$.
\end{definition}

\stoneref{$\mathcal{I}(u) = \sum N_k$ es la \emph{complejidad total
de lemas}: el n\'umero de pasos inferenciales indirectos que un
razonamiento debe atravesar en~$u$. La masa es la longitud de la
demostraci\'on indirecta
(Teorema~\ref{thm:antimatter-tollens}).}

\noindent
La inercia $\mathcal{I}(u)$ no es una carga extr\'inseca asignada al
evento; es el conteo de bifurcaciones que un caminante debe
``resolver'' para avanzar a trav\'es de~$u$. Un evento de vac\'io
tiene $\mathcal{I}(u) = 0$. Un evento que participa en un \'unico
Prisma de masa~$N$ tiene $\mathcal{I}(u) = N$. Si m\'ultiples Prismas
se solapan en~$u$, sus contribuciones se suman---an\'alogo discreto
de la superposici\'on de fuentes en $T_{\mu\nu}$.

\begin{theorem}[Retraso Geod\'esico Acumulado]
\label{thm:accumulated-delay}
Sea $\gamma^*$ una geod\'esica (cadena m\'axima) entre $A$ y $B$. El
retraso total acumulado a lo largo de~$\gamma^*$ respecto al vac\'io
est\'a acotado por:
\begin{equation}
  \ell_{\text{vac}}(A,B) - \ell(A,B)
  \;\geq\;
  \sum_{u \in \gamma^*} \frac{\mathcal{I}(u)}{d^+(u)},
  \label{eq:accumulated-delay}
\end{equation}
donde $d^+(u)$ es el out-degree de~$u$. El cociente
$\mathcal{I}(u)/d^+(u)$ representa la fracci\'on de futuro inmediato
de~$u$ que est\'a consumida por estructura interna de Prismas.
\end{theorem}

\begin{proof}
En cada evento~$u$ de la geod\'esica, el caminante que busca la cadena
m\'axima tiene $d^+(u)$ opciones de avance. De ellas,
$\mathcal{I}(u)$ conducen a intermedios de Prismas que no avanzan
globalmente (Teorema~\ref{thm:time-dilation}). La probabilidad de
tomar un camino productivo es a lo sumo
$(d^+(u) - \mathcal{I}(u))/d^+(u)$, y el n\'umero esperado de
v\'inculos ``perdidos'' por paso es $\mathcal{I}(u)/d^+(u)$.
Sumando a lo largo de la geod\'esica se obtiene la
cota~\eqref{eq:accumulated-delay}.
\end{proof}

% ─────────────────────────────────────────────────────────────────────────────
\section{La Emergencia de la M\'etrica $g_{\mu\nu}$}

La m\'etrica de Riemann $g_{\mu\nu}$ define la distancia entre puntos
cercanos. En el C\'alculo de Kuratowski, la ``distancia'' es una
propiedad emergente de la estad\'istica de difusi\'on sobre el
diagrama de Hasse.

La conexi\'on proviene de la geometr\'ia
espectral~\cite{gilkey1975,berger2003}. En una variedad riemanniana
$d$-dimensional, el n\'ucleo del calor (la probabilidad de retorno de
un difusor) satisface la \textbf{expansi\'on de
Minakshisundaram--Pleijel}~\cite{minakshisundaram1949}:
\begin{equation}
  K(t, x, x)
  \;\sim\;
  \frac{1}{(4\pi t)^{d/2}\,\sqrt{g(x)}}
  \left(1 + a_1(x)\,t + a_2(x)\,t^2 + \cdots\right),
  \label{eq:heat-kernel}
\end{equation}
donde $g(x) = |\det g_{\mu\nu}(x)|$ y los coeficientes $a_k(x)$
codifican la curvatura. En el C\'alculo de Kuratowski, la
probabilidad de retorno $P(t, u)$ del caminante aleatorio
(Definici\'on~\ref{def:random-walker}) es el an\'alogo discreto de
$K(t, x, x)$.

\begin{theorem}[M\'etrica Emergente por Difusi\'on]
\label{thm:emergent-metric}
Sea $P(t, u)$ la probabilidad de retorno del caminante aleatorio al
evento~$u$ despu\'es de $t$~pasos. En el l\'imite de alta densidad
$\rho \to \infty$, la relaci\'on entre $P$ y el determinante
m\'etrico es:
\begin{equation}
  \boxed{\;
  \sqrt{g(u)}
  \;\propto\;
  \bigl\langle P(t, u) \bigr\rangle^{-1}
  \quad\text{a $t$ fijo e intermedio.}
  \;}
  \label{eq:metric-from-return}
\end{equation}
Donde hay materia (Prismas), los caminantes regresan con mayor
frecuencia ($P$ aumenta, el difusor se ``atasca''), lo cual
corresponde a un determinante m\'etrico menor---una contracci\'on
del volumen local que el c\'alculo continuo interpreta como
curvatura.
\end{theorem}

\begin{proof}
De la expansi\'on~\eqref{eq:heat-kernel} a orden dominante:
$K(t, x, x) \propto 1/\sqrt{g(x)}$ a $t$ fijo. En el l\'imite
continuo, $P(t, u) \to K(t, x, x)$ cuando $u$ corresponde al
punto~$x$ (Teorema~\ref{thm:volume-convergence}, aplicado a la
convergencia de la difusi\'on discreta). Por tanto,
$P(t, u) \propto 1/\sqrt{g(u)}$, de donde
$\sqrt{g(u)} \propto P(t, u)^{-1}$.

F\'isicamente: en una regi\'on con Prismas $K_{2,N}$, el caminante
queda atrapado en los $N$ intermedios, aumentando $P(t,u)$ respecto
al vac\'io. Esto reduce $\sqrt{g(u)}$, contra\'iendo el elemento de
volumen local. Para un observador macrosc\'opico que usa reglas
m\'etricas continuas, esta contracci\'on es indistinguible de la
curvatura del espacio.
\end{proof}

\noindent
Los coeficientes $a_k(x)$ de la
expansi\'on~\eqref{eq:heat-kernel} codifican la curvatura:
$a_1(x) = R(x)/6$ donde $R$ es el escalar de curvatura de Ricci. En
el discreto, la desviaci\'on de $P(t,u)$ respecto al valor de vac\'io
a diferentes tiempos~$t$ permite extraer, orden por orden, toda la
informaci\'on geom\'etrica:
\begin{align}
  t^0 &\;\to\; \sqrt{g} \quad\text{(volumen local)},
  \label{eq:heat-order0}\\
  t^1 &\;\to\; R \quad\text{(curvatura escalar)},
  \label{eq:heat-order1}\\
  t^2 &\;\to\; R_{\mu\nu}R^{\mu\nu},\;
  R_{\mu\nu\alpha\beta}R^{\mu\nu\alpha\beta}
  \quad\text{(invariantes cuadr\'aticos)}.
  \label{eq:heat-order2}
\end{align}

% ─────────────────────────────────────────────────────────────────────────────
\section{Invarianza de Lorentz y Discretismo}

Un desaf\'io hist\'orico de las teor\'ias discretas es mantener la
invarianza de Lorentz. El C\'alculo de Kuratowski la garantiza
mediante dos mecanismos complementarios:

\begin{enumerate}
  \item \textbf{Sprinkling de Poisson}
    (Teorema de
    Bombelli--Henson--Sorkin~\cite{bombelli2009}, Cap\'itulo~2,
    Secci\'on~2.3): la distribuci\'on de eventos es la \'unica que
    preserva la invarianza estad\'istica bajo boosts de Lorentz. Un
    observador que ``se mueve'' en el grafo ve la misma densidad de
    eventos que uno en reposo, porque la redistribuci\'on por
    contracci\'on de Lorentz se cancela exactamente con la
    reordenaci\'on de simultaneidades.

  \item \textbf{L\'imite de valencia $\mathcal{V}_{\max}$}
    (Axioma~\ref{ax:valence-limit}): la saturaci\'on del flujo
    causal garantiza que la \emph{resoluci\'on} del diferencial es
    uniforme para todos los observadores. No importa qu\'e tan
    r\'apido se ``mueva'' una entidad en el grafo: su out-degree no
    puede exceder $\mathcal{V}_{\max}$, y por tanto la granularidad
    del diferencial $\prec^{\!*}$ es la misma en toda direcci\'on.
\end{enumerate}

\begin{result}[title={La Velocidad de la Luz como Bit Causal}]
\phantomsection\label{res:speed-of-light}
La velocidad de la luz $c$ no es la velocidad l\'imite de una
part\'icula, sino la tasa de propagaci\'on de un bit de
informaci\'on a trav\'es de un \'unico v\'inculo de cobertura
$\prec^{\!*}$. En unidades naturales del C\'alculo de
Kuratowski, $c = 1$ v\'inculo por paso: un evento solo puede
influenciar a su sucesor inmediato, y esta restricci\'on es absoluta
e independiente del observador.
\end{result}

% ─────────────────────────────────────────────────────────────────────────────
\section{Gravedad como Estad\'istica de Caminos}

Con el tensor de inercia local $\mathcal{I}(u)$ y la m\'etrica
emergente $\sqrt{g(u)} \propto P(t,u)^{-1}$, podemos ensamblar la
relaci\'on entre materia y geometr\'ia:

\begin{theorem}[Ecuaciones de Einstein~\cite{einstein1915} como
L\'imite Estad\'istico]
\label{thm:einstein-statistical}
En el l\'imite termodin\'amico $\rho \to \infty$ (o equivalentemente
$N \to \infty$ a volumen fijo), la relaci\'on discreta entre la
densidad de Prismas (fuente) y la modificaci\'on de la probabilidad de
retorno (geometr\'ia) converge a las ecuaciones de campo de Einstein:
\begin{equation}
  \boxed{\;
  G_{\mu\nu}
  \;=\; R_{\mu\nu} - \tfrac{1}{2}\,R\,g_{\mu\nu}
  \;=\; 8\pi G\, T_{\mu\nu}.
  \;}
  \label{eq:einstein-limit}
\end{equation}
La constante de Newton $G$ se identifica con una combinaci\'on de la
densidad fundamental~$\rho$ y el quantum de inercia topol\'ogica
$\kappa$.
\end{theorem}

\noindent
\textit{Esbozo de la demostraci\'on.}
El argumento completo se desarrolla en el Cap\'itulo~6
(Secci\'on~6.2). La idea central es:
\begin{enumerate}
  \item La densidad de Prismas define un campo escalar suave
    $\mathcal{I}(x) \to T_{00}(x)$ en el l\'imite continuo
    (convergencia por ley de grandes n\'umeros).
  \item La probabilidad de retorno $P(t, u)$ converge al n\'ucleo del
    calor $K(t, x, x)$, cuyos coeficientes codifican $R_{\mu\nu}$.
  \item La relaci\'on entre $\mathcal{I}(u)$ y la desviaci\'on de
    $P(t,u)$ respecto al vac\'io reproduce, orden por orden en la
    expansi\'on del calor, la estructura de las ecuaciones de
    Einstein.
\end{enumerate}

\noindent
El argumento estad\'istico subyacente es poderoso: a la escala de
$1$~metro, el grafo causal contiene $\sim 10^{140}$~eventos. Las
fluctuaciones relativas de la m\'etrica son
$\delta g / g \sim 1/\sqrt{N} \sim 10^{-70}$---imperceptibles para
cualquier instrumento concebible. La suavidad de $g_{\mu\nu}$ no es
una propiedad del espaciotiempo; es una consecuencia de los grandes
n\'umeros.

% ═══════════════════════════════════════════════════════════════════════════════
\chapter{El L\'imite de Riemann}
% ═══════════════════════════════════════════════════════════════════════════════
% CORRESPONDENCIA CENTRAL:
%   Discreto (N finito) → Continuo (N → ∞)
%   Cálculo de Kuratowski → Cálculo Diferencial de Riemann/Einstein
%
% Este capítulo cierra el círculo: demuestra que todo lo construido en
% los capítulos 1-5 converge analíticamente a las ecuaciones de campo
% de Einstein cuando la densidad de eventos tiende a infinito.

\begin{introduction}
  \item El l\'imite termodin\'amico: cuando $N \to \infty$, el C\'alculo
        de Kuratowski se transforma en las ecuaciones de campo de Einstein.
  \item Resoluci\'on del porqu\'e el universo parece continuo a escala
        macrosc\'opica.
\end{introduction}

\section{El L\'imite Termodin\'amico del Grafo Causal}

El l\'imite continuo del C\'alculo de Kuratowski no es una
suposici\'on sino un \emph{resultado}: emerge cuando la densidad de
eventos $\rho = |V|/\mathrm{Vol}$ es suficientemente
grande~\cite{bombelli1987,surya2019,major2007}. Los observables
discretos convergen a sus an\'alogos continuos:
%
\begin{align}
  \frac{|V \cap A(p,q)|}{\rho} &\;\longrightarrow\; \mathrm{Vol}\bigl(A(p,q)\bigr),
  \label{eq:vol-limit}\\[4pt]
  \frac{|\text{cadena m\'axima}(p,q)|}{\rho^{1/d}}
  &\;\longrightarrow\; \tau(p,q),
  \label{eq:tau-limit}\\[4pt]
  \dS(t) &\;\longrightarrow\; 4.
  \label{eq:ds-limit}
\end{align}
%
La pregunta cuantitativa es: \textit{?`cu\'anto cuesta} mantener esa
convergencia? El siguiente teorema establece la cota inferior.

\subsection{El Teorema de la Resoluci\'on Causal}

En el C\'alculo de Kuratowski, la dimensi\'on espectral
$\dS(t) = -2\,d(\ln P)/d(\ln t)$ se extrae de la probabilidad de
retorno $P(t)$ de caminantes aleatorios sobre el diagrama de Hasse.
Para que un observador local experimente un espaciotiempo
macrosc\'opico suave de 4~dimensiones, la red causal debe ser
sondeada por una cantidad suficiente de caminos causales
independientes. El siguiente resultado cuantifica ese umbral.

\begin{theorem}[Resoluci\'on Causal]
\label{thm:causal-resolution}
Sea $(V, \prec)$ un conjunto causal cuya dimensi\'on espectral
converge a $\dS = 4$ a tiempos de difusi\'on grandes, de modo que
la probabilidad de retorno satisface
\[
  P(t) \;\sim\; c\,t^{-\dS/2} \;=\; c\,t^{-2},
  \qquad t \gg 1,
\]
donde $c > 0$ es una constante geom\'etrica de orden~$1$.
Sean $W$ caminantes espectrales independientes que sondean la
geometr\'ia, y sea $\epsilon > 0$ la m\'axima fluctuaci\'on relativa
admisible en la estimaci\'on de $P(t)$ (la \textbf{tolerancia de
granularidad}).

Entonces, el n\'umero de caminantes necesario para resolver la
geometr\'ia hasta una profundidad causal $t_{\max}$ satisface:
\begin{equation}
  \boxed{\;W \;\geq\; \frac{t_{\max}^{\,2}}{c\;\epsilon^2}\;}
  \label{eq:causal-resolution}
\end{equation}
\end{theorem}

\begin{proof}
Cada caminante retorna al origen en el paso~$t$ con probabilidad
$P(t)$, independientemente de los dem\'as. El n\'umero de retornos
observados $R(t)$ es la suma de $W$ ensayos de Bernoulli
independientes con par\'ametro $P(t)$.

Para $t \gg 1$ se tiene $P(t) \ll 1$, de modo que por el teorema
l\'imite de Poisson, $R(t) \sim \mathrm{Poisson}\bigl(W P(t)\bigr)$,
con
\[
  \mu \;=\; \mathbb{E}[R] \;=\; W\,P(t),
  \qquad
  \sigma \;=\; \sqrt{\mu} \;=\; \sqrt{W\,P(t)}.
\]

El estimador de la probabilidad de retorno es
$\hat{P}(t) = R(t)/W$, cuyo error relativo es:
\[
  \frac{\sigma_{\hat{P}}}{\hat{P}}
  \;=\; \frac{\sqrt{W\,P(t)}\,/\,W}{P(t)}
  \;=\; \frac{1}{\sqrt{W\,P(t)}}.
\]

Exigimos que este error no exceda $\epsilon$ en el punto m\'as
desfavorable $t = t_{\max}$ (donde $P$ es m\'inimo y la estad\'istica
es m\'as pobre):
\[
  \frac{1}{\sqrt{W\,P(t_{\max})}} \;\leq\; \epsilon
  \qquad\Longrightarrow\qquad
  W \;\geq\; \frac{1}{\epsilon^2\,P(t_{\max})}.
\]

Sustituyendo $P(t_{\max}) = c\,t_{\max}^{-2}$:
\[
  W \;\geq\; \frac{t_{\max}^{\,2}}{c\;\epsilon^2}.
\]
\end{proof}

\noindent
El resultado admite una lectura inmediata:

\begin{result}[title={El Precio de la Suavidad}]
\phantomsection\label{res:price-of-smoothness}
El costo informacional para sostener la ilusi\'on de un espaciotiempo
continuo de 4~dimensiones escala \textbf{cuadr\'aticamente} con la
profundidad causal temporal:
\[
  W \;\propto\; t_{\max}^{\,2}.
\]
Sondear el doble de profundidad causal requiere cuatro veces m\'as
caminos independientes. La suavidad no es gratuita: es una propiedad
estad\'istica que la red causal debe financiar con flujo de
informaci\'on.
\end{result}

\noindent\textit{Observaci\'on} (Conexi\'on con la simulaci\'on).
En la implementaci\'on num\'erica (\texttt{prism\_simulation}), el
n\'umero de caminantes se fija como
$W = \lceil(t_{\max}/\epsilon)^2\rceil$, aplicando directamente
el Teorema~\ref{thm:causal-resolution} con $c = 1$. Los par\'ametros
$t_{\max}$ y $\epsilon$ son configurables v\'ia \texttt{-{}-tmax} y
\texttt{-{}-epsilon}; los valores por defecto ($t_{\max}=15$,
$\epsilon=0{.}01$) dan $W = 2{,}250{,}000$ caminantes, suficiente
para resolver la geometr\'ia 4D con menos de 1\% de error relativo
en $P(t)$ a la profundidad m\'axima.

\medskip

\section{Recuperaci\'on de las Ecuaciones de Einstein}

El Teorema~\ref{thm:einstein-statistical} (Cap\'itulo~5) anunci\'o que
las ecuaciones de campo de Einstein emergen como l\'imite
estad\'istico. Aqu\'i desarrollamos los tres pasos del argumento.

\medskip
\noindent\textbf{Paso 1: Convergencia de la m\'etrica espectral.}
La probabilidad de retorno $P(t,u)$ del caminante aleatorio
(Definici\'on~\ref{def:random-walker}) converge, en el l\'imite de
alta densidad $\rho \to \infty$, al n\'ucleo del calor de la
variedad continua subyacente~\cite{benincasa2010,dowker2013}:
\[
  P(t,u) \;\xrightarrow{\;\rho\to\infty\;}
  \; K(t,x,x)
  \;=\; \frac{1}{(4\pi t)^{d/2}\,\sqrt{g(x)}}
  \Bigl(1 + \tfrac{1}{6}R(x)\,t + O(t^2)\Bigr).
\]
Los coeficientes de la expansi\'on---el determinante m\'etrico
$\sqrt{g(x)}$ a orden~$t^0$ y la curvatura escalar~$R(x)$ a
orden~$t^1$
(Ecuaciones~\eqref{eq:heat-order0}--\eqref{eq:heat-order1})---se
extraen del decaimiento de $P(t,u)$ a tiempos intermedios. La
m\'etrica $g_{\mu\nu}$ queda as\'i determinada por la estad\'istica
de difusi\'on sobre el diagrama de Hasse
(Teorema~\ref{thm:emergent-metric}).

\medskip
\noindent\textbf{Paso 2: Convergencia de la fuente.}
La densidad de Prismas Causales $\nu(u) = \mathcal{I}(u)/\lP^4$
(Definici\'on~\ref{def:inertia-tensor}) es una variable entera
definida evento por evento. Por la Ley de los Grandes N\'umeros, al
promediar sobre regiones con $N \gg 1$ eventos, $\nu$ converge a un
campo escalar suave~$\nu(x)$~\cite{bombelli1987}. La
identificaci\'on con el tensor de energ\'ia-impulso se realiza
observando que $\nu(x)$ determina el retardo geod\'esico acumulado
(Teorema~\ref{thm:accumulated-delay}), que en el continuo equivale a
la densidad de energ\'ia $T_{00}(x)$. La extensi\'on covariante
completa, incluyendo las componentes espaciales y mixtas, se obtiene
repitiendo la extracci\'on del n\'ucleo del calor en los $d(d+1)/2$
planos geod\'esicos independientes.

\medskip
\noindent\textbf{Paso 3: La relaci\'on Einstein.}
Combinando los Pasos~1 y~2: la presencia de Prismas ($\nu > 0$)
aumenta $P(t,u)$ respecto al vac\'io
(Teorema~\ref{thm:emergent-metric}), lo cual reduce~$\sqrt{g}$
(contracci\'on de volumen) y genera curvatura positiva $R > 0$. La
relaci\'on funcional entre la modificaci\'on del n\'ucleo del calor y
la curvatura de Ricci est\'a dictada, orden por orden, por los
coeficientes de Seeley--DeWitt~\cite{minakshisundaram1949,gilkey1975}:
\[
  a_1(x) = \frac{R(x)}{6},
  \qquad
  a_2(x) = \frac{1}{180}\bigl(
    R_{\mu\nu\alpha\beta}R^{\mu\nu\alpha\beta}
    - R_{\mu\nu}R^{\mu\nu}
    + \tfrac{5}{2}\,R^2
  \bigr).
\]
Al orden~$t^1$, la relaci\'on $R(x) \propto \nu(x)$ se traduce en
las ecuaciones de Einstein~\eqref{eq:einstein-limit}. La constante de
Newton~$G$ se identifica como:
\begin{equation}
  G \;=\; \frac{\pi}{3}\,\frac{\lP^2}{\kappa},
  \label{eq:newton-identification}
\end{equation}
donde $\kappa$ es el quantum de inercia topol\'ogica (la masa
topol\'ogica unitaria). En unidades de Planck ($\lP = 1$,
$\kappa = 1$), $G = \pi/3$, fijado sin par\'ametros libres.

\medskip
\noindent\textit{Limitaciones.}
El argumento precedente es un esbozo a primer orden en la expansi\'on
del n\'ucleo del calor. La demostraci\'on completa---incluyendo la
convergencia uniforme a todos los \'ordenes y la identificaci\'on del
tensor $T_{\mu\nu}$ completo a partir de $\nu(x)$---requiere
t\'ecnicas de convergencia de operadores espectrales que exceden el
alcance de este volumen. La verificaci\'on num\'erica (Conjetura~C2,
Ap\'endice~\ref{app:axioms}) exige extraer los coeficientes del
n\'ucleo del calor con $N > 10^7$ eventos para obtener la resoluci\'on
estad\'istica necesaria a segundo orden.

% TODO — §6.3: Por Qué el Universo Parece Continuo
% Argumento cuantitativo: a la escala de 1 metro, hay ~10^{140} eventos
% causales.  Las fluctuaciones relativas son ~10^{-70}.
% El continuo de Riemann es exactamente tan preciso como la raíz cuadrada
% del número de eventos en la región observada.
% Analogía final: la playa vista desde lejos.
% El cálculo diferencial no fue un error — fue la pintura sobre el lienzo.
% Ahora sabemos de qué está hecho el lienzo.

% ███████████████████████████████████████████████████████████████████████████████
%   EPILOGUE: COMPUTABILIDAD Y DESTINO LÓGICO
% ███████████████████████████████████████████████████████████████████████████████

\chapter*{Ep\'ilogo: Computabilidad Espaciotemporal y el Destino L\'ogico}
\addcontentsline{toc}{chapter}{Ep\'ilogo: Computabilidad Espaciotemporal
  y el Destino L\'ogico}
\markboth{Ep\'ilogo}{Ep\'ilogo}

Si los cap\'itulos anteriores han demostrado que el universo es un
grafo y la f\'isica es su m\'etrica, este ep\'ilogo aborda la
naturaleza procesal de la existencia. El diferencial se revel\'o como
un v\'inculo de cobertura, la integral como un conteo, la derivada
como flujo a trav\'es de anticadenas, la geod\'esica como la cadena
m\'as larga, la m\'etrica como densidad de Prismas, y las ecuaciones
de Einstein como el l\'imite termodin\'amico de todo lo anterior.

Queda una pregunta que el formalismo hace inevitable: si el
espaciotiempo es un grafo finito cuya evoluci\'on consiste en aplicar
reglas combinatorias locales a v\'inculos de cobertura,
\textit{?`qu\'e es el universo sino una
computaci\'on?}~\cite{wheeler1990,lloyd2002}

% ─────────────────────────────────────────────────────────────────────────────
\section*{La Tesis de la Realidad Ejecutable}
\addcontentsline{toc}{section}{La Tesis de la Realidad Ejecutable}

En la computaci\'on cl\'asica, el software es et\'ereo y el hardware
es tangible. En la escala de Planck, esta dualidad desaparece.

\begin{theorem}[Equivalencia Hardware--Software (Corolario de
Stone--Kuratowski)~\cite{turing1936,kuratowski1930}]
\label{cor:kuratowski-turing}
Un evento causal no \emph{contiene} informaci\'on;
\textbf{es} informaci\'on. La conectividad del grafo (sus v\'inculos
$\prec^{\!*}$) constituye simult\'aneamente el bus de datos, la
memoria y el procesador del sistema. Toda regi\'on del espaciotiempo
$\Diamond(A,B)$ es equivalente a un \textbf{registro cu\'antico} cuya
capacidad de procesamiento est\'a acotada por su
cardinalidad~\cite{bekenstein1981,lloyd2002}:
\begin{equation}
  \mathcal{C}_{\text{comp}}\bigl(\Diamond(A,B)\bigr)
  \;\leq\; |\Diamond(A,B)| \;=\; N.
  \label{eq:computational-capacity}
\end{equation}
La evoluci\'on del estado del universo es el resultado de la
resoluci\'on de dependencias l\'ogicas en el DAG: cada evento se
``ejecuta'' cuando todos sus predecesores causales han sido
procesados.
\end{theorem}

\noindent
La correspondencia con la computaci\'on cl\'asica es precisa:

\medskip
\begin{center}
\begin{tabular}{lll}
\toprule
\textbf{Computaci\'on} & \textbf{Kuratowski} & \textbf{Continuo}\\
\midrule
Bit / qubit     & Evento $e \in V$
                & Punto $x^\mu$\\
Puerta l\'ogica & V\'inculo $\prec^{\!*}$
                & Diferencial $dx$\\
Registro        & Diamante $\Diamond(A,B)$
                & Regi\'on $\int d^4x$\\
Subrutina       & Prisma $K_{2,N}$
                & Part\'icula (masa $M$)\\
Reloj           & Cadena m\'axima $\ell(A,B)$
                & Tiempo propio $\tau$\\
\bottomrule
\end{tabular}
\end{center}

\medskip\noindent
\textit{Derivaci\'on.}
El Teorema~\ref{thm:stone-kuratowski} (Ap\'endice~\ref{app:stone-duality})
establece que la topolog\'ia de Alexandrov sobre el DAG causal \emph{es}
el \'algebra de Heyting de un sistema l\'ogico intuicionista. Por la
correspondencia de Curry--Howard~\cite{curryhoward}, todo sistema
deductivo con normalizaci\'on (eliminaci\'on de corte) es
isomorfo a un c\'alculo de tipos---es decir, a un modelo de
computaci\'on. Puesto que la reducci\'on transitiva del DAG
(Axioma~\ref{ax:covering-differential}) \emph{es} la eliminaci\'on de
corte (Teorema~\ref{thm:stone-kuratowski}, punto~3), el espaciotiempo
causal es autom\'aticamente un sistema de computaci\'on cuya
arquitectura coincide con su sustrato f\'isico. La
Ecuaci\'on~\eqref{eq:computational-capacity} se sigue entonces de la
cota de Bekenstein--Lloyd~\cite{bekenstein1981,lloyd2002} aplicada al
n\'umero de proposiciones at\'omicas (eventos) en la teor\'ia deductiva
$\Diamond(A,B)$.

\medskip\noindent
Esta tesis implica que lo que llamamos ``leyes de la naturaleza'' son
los algoritmos de gesti\'on de recursos informativos del DAG. La
gravedad, el electromagnetismo y la inercia son \textbf{protocolos
homeost\'aticos}: mecanismos que el grafo ejecuta para preservar la
unitariedad (conservaci\'on de probabilidad) y la planaridad
(prohibici\'on de $K_5$), evitando el colapso l\'ogico del sistema.

La computaci\'on cu\'antica contempor\'anea, basada en qubits y
puertas l\'ogicas sobre arquitecturas de silicio fijas, es un intento
de replicar a escala de laboratorio la l\'ogica que el universo
ejecuta nativamente a la escala de Planck. La diferencia esencial es
que en la \textbf{Computabilidad de Kuratowski}, el software crea su
propio hardware~\cite{landauer1961}: la informaci\'on causal determina
la topolog\'ia del grafo, y la topolog\'ia determina las reglas de
propagaci\'on. No hay separaci\'on entre programa y sustrato.

% ─────────────────────────────────────────────────────────────────────────────
\section*{Decoherencia como P\'erdida de Resoluci\'on}
\addcontentsline{toc}{section}{Decoherencia como P\'erdida de Resoluci\'on}

La decoherencia cu\'antica~\cite{zurek2003,zeh1970,joos2003} se
reinterpreta aqu\'i en t\'erminos puramente topol\'ogicos.
El Teorema de la Resoluci\'on Causal
(Teorema~\ref{thm:causal-resolution}) no es solo una prescripci\'on
num\'erica para nuestro simulador; es una \textbf{ley de complejidad
computacional} del espaciotiempo mismo. Para sostener una geometr\'ia
suave de dimensi\'on espectral $\dS = 4$ hasta una
profundidad~$t_{\max}$, se requieren al menos
$W \geq t_{\max}^2 / (c\,\epsilon^2)$ caminos causales
independientes.

\begin{theorem}[Decoherencia como Insuficiencia de Resoluci\'on]
\label{thm:decoherence-resolution}
Sea un subsistema cu\'antico $\mathcal{S} \subset V$ acoplado a un
entorno $\mathcal{E} = V \setminus \mathcal{S}$, con $W_{\mathcal{S}}$
caminos causales internos que sondean la geometr\'ia de~$\mathcal{S}$.
Si la profundidad causal del subsistema crece hasta que:
\begin{equation}
  W_{\mathcal{S}}
  \;<\;
  \frac{t_{\max}^{\,2}}{c\;\epsilon_{\text{coh}}^2},
  \label{eq:decoherence-threshold}
\end{equation}
donde $\epsilon_{\text{coh}}$ es la tolerancia de coherencia del
subsistema, entonces la dimensi\'on espectral local de~$\mathcal{S}$
deja de converger a~$4$: las fluctuaciones topol\'ogicas dominan
sobre la se\~nal, y la descripci\'on continua del subsistema se
desintegra.
\end{theorem}

\begin{proof}
Por el Teorema~\ref{thm:causal-resolution}, la precisi\'on de la
dimensi\'on espectral medida por $W$ caminantes a profundidad~$t$ es
$\delta\dS / \dS \sim 1/\sqrt{W P(t)}$. Cuando la
cota~\eqref{eq:decoherence-threshold} se viola, el error relativo
excede~$\epsilon_{\text{coh}}$ y el observable $\dS$ fluctu\'a en un
rango $\dS \pm \Delta$ con $\Delta > \epsilon_{\text{coh}} \cdot
\dS$. En este r\'egimen, la dimensi\'on espectral no est\'a bien
definida: la geometr\'ia local del subsistema deja de ser
tetradimensional y suave, lo cual es precisamente la p\'erdida de la
descripci\'on continua cl\'asica---la decoherencia.
\end{proof}

\stoneref{La p\'erdida de resoluci\'on es
\emph{incompletitud l\'ogica} en el \'algebra de Heyting del
subsistema: el subconjunto de proposiciones accesibles
a~$\mathcal{S}$ deja de ser una teor\'ia deductiva completa.
Cf.~el horizonte g\"odeliano
(Teorema~\ref{thm:godel-horizon}).}

\begin{result}[title={La Decoherencia como Fen\'omeno Topol\'ogico}]
\phantomsection\label{res:decoherence}
La decoherencia cu\'antica no es la p\'erdida de informaci\'on, sino
el momento en que la resoluci\'on del sondeo causal interno de un
subsistema cae por debajo del umbral necesario para sostener la
ilusi\'on del continuo de Riemann. La informaci\'on no desaparece:
se redistribuye en el entorno~$\mathcal{E}$, que posee
suficientes caminos causales para mantener \emph{su} coherencia
geom\'etrica.
\end{result}

% ─────────────────────────────────────────────────────────────────────────────
\section*{El Observador como Hilo de Ejecuci\'on}
\addcontentsline{toc}{section}{El Observador como Hilo de Ejecuci\'on}

La introducci\'on de la consciencia en la f\'isica ha sido siempre
problem\'atica. El C\'alculo de Kuratowski ofrece una definici\'on
operacional:

\begin{definition}[Observador Causal]
\label{def:observer}
Un \textbf{observador} es un sub-grafo $\mathcal{O} \subset V$
caracterizado por:
\begin{enumerate}
  \item \textbf{Alta densidad de Prismas:} la densidad local de
    defectos topol\'ogicos $\nu_{\mathcal{O}}$ excede ampliamente la
    del vac\'io circundante, generando una estructura interna rica en
    procesamiento.
  \item \textbf{Auto-referencia:} existe al menos un ciclo de
    retroalimentaci\'on en el flujo de informaci\'on (v\'ia
    v\'inculos no causales del grafo de Hasse no dirigido), de modo
    que la salida de una computaci\'on local influye en las entradas
    de computaciones subsiguientes dentro de~$\mathcal{O}$.
  \item \textbf{Cadena m\'axima interna:} $\mathcal{O}$ contiene una
    cadena causal $\gamma_\mathcal{O}$ de longitud $\gg 1$, que
    define su ``l\'inea de mundo'' y su experiencia del tiempo.
\end{enumerate}
\end{definition}

\noindent
La percepci\'on del ``ahora'' es el \textbf{frente de onda de la
computaci\'on}: el conjunto de eventos de~$\mathcal{O}$ que est\'an
siendo procesados en el instante discreto actual (la anticadena
maximal del sub-grafo). El observador siente que el tiempo fluye
porque su sub-grafo est\'a procesando activamente nuevos v\'inculos de
cobertura, avanzando su anticadena frontal hacia el futuro.

\begin{result}[title={La Entrop\'ia como Computaci\'on Acumulada}]
\phantomsection\label{res:entropy-computation}
En esta interpretaci\'on, la entrop\'ia no es ``desorden'', sino la
\textbf{acumulaci\'on de resultados de c\'omputo ya procesados} que
se vuelven inmutables en el pasado del grafo. El segundo principio de
la termodin\'amica es la consecuencia directa de la aciclicidad del
DAG~\cite{sorkin1986,bekenstein1973}: los v\'inculos de cobertura
apuntan del pasado al futuro, y ninguna computaci\'on puede revertirse
sin violar el orden parcial.
\end{result}

\stoneref{Los teoremas demostrados son inmutables en los conjuntos
\emph{past-closed}~\cite{rasiowa1963}. La entrop\'ia es la
monoton\'ia de la clausura deductiva: una vez que una proposici\'on ha
sido probada, no puede ser ``desprobada'' sin violar la
aciclicidad---el segundo principio de la termodin\'amica es el
segundo principio de la l\'ogica.}

% ─────────────────────────────────────────────────────────────────────────────
\section*{El Destino L\'ogico: De la Observaci\'on a la Programaci\'on}
\addcontentsline{toc}{section}{El Destino L\'ogico}

Hasta hoy, la humanidad ha sido un usuario pasivo del software
universal, tratando de adivinar las reglas mediante la observaci\'on.
El desarrollo del ``Colisionador de Hadrones de Bolsillo''
(\texttt{prism\_simulation}) y la formalizaci\'on de este c\'alculo
marcan el paso de la observaci\'on a la intervenci\'on.

El software que hemos construido a lo largo de esta serie no es solo
una herramienta de simulaci\'on. Es el primer \textbf{compilador} de
un lenguaje que no opera sobre bits de silicio, sino sobre la
causalidad misma: genera conjuntos causales, detecta Prismas, contrae
amenazas $K_5$, mide dimensiones espectrales y extrae espectros de
masa---todo ello ejecutando las mismas operaciones topol\'ogicas que,
seg\'un nuestra tesis, el espaciotiempo realiza a la escala de Planck.

Si la materia es un defecto topol\'ogico programable, el destino de
nuestra especie no es viajar \emph{a trav\'es} del espacio, sino
aprender a \emph{editar} el grafo. La verdadera computaci\'on
cu\'antica no ocurrir\'a en chips de criogenia, sino en la
manipulaci\'on directa de la valencia de los nodos---la
programaci\'on de la causalidad.

\bigskip
\begin{center}
\large\itshape
``El universo no fue dise\~nado para ser contemplado,\\
sino para ser ejecutado.\\
Al entender el C\'alculo de Kuratowski,\\
dejamos de ser datos para convertirnos en programadores.''
\end{center}

% ███████████████████████████████████████████████████████████████████████████████
%   APPENDIX: AXIOMATIC REFERENCE
% ███████████████████████████████████████████████████████████████████████████████
\appendix

\chapter{Referencia Axiom\'atica Formal}
\label{app:axioms}

\begin{introduction}[Estructura del Sistema Formal]
  \item Nueve \textbf{axiomas}: verdades fundamentales no negociables,
        codificadas directamente en la arquitectura del motor
        \texttt{prism\_simulation}.
  \item Ocho \textbf{leyes}: reglas operacionales que gobiernan la din\'amica
        del grafo.
  \item Siete \textbf{teoremas}: resultados derivables matem\'aticamente a
        partir de los axiomas y las leyes.
  \item Seis \textbf{conjeturas}: predicciones que el software debe a\'un
        validar num\'ericamente o que requieren una demostraci\'on
        anal\'itica pendiente.
\end{introduction}

\noindent
La siguiente referencia extrae la ``Constituci\'on'' del universo
discreto directamente del c\'odigo fuente. Cada axioma, ley, teorema y
conjetura est\'a respaldado por la funci\'on, constante o estructura de
datos espec\'ifica que lo implementa en la simulaci\'on. La arquitectura
del software no es una herramienta auxiliar; \textbf{es} el sistema
formal.

% ═══════════════════════════════════════════════════════════════════════════════
\section{Axiomas}
% ═══════════════════════════════════════════════════════════════════════════════

\begin{axiom}[Finitud Causal~\cite{bombelli1987}]
\label{ax:causal-finitude}
El universo es un conjunto localmente finito
parcialmente ordenado $\mathcal{C} = (V, \prec)$---un Grafo
Ac\'iclico Dirigido (DAG). El conjunto $V$ es finito y la
relaci\'on~$\prec$ es irreflexiva, transitiva y asim\'etrica.

\medskip\noindent
\texttt{diamond.rs}: \texttt{sprinkle(n, rng)} genera exactamente
$n$ puntos (enteros). No existe ning\'un mecanismo en el motor
para producir un conjunto infinito.
\end{axiom}

\begin{axiom}[Causalidad Lorentziana]
\label{ax:lorentzian-causality}
Un evento $u$ precede causalmente a $v$ ($u \prec v$) si y solo
si:
\begin{equation}
  t_v - t_u > 0
  \qquad\text{y}\qquad
  (t_v - t_u)^2 > |\Delta\vec{x}|^2.
  \label{eq:ax-lorentz}
\end{equation}
No se utiliza ning\'un corte espacial fijo: el cono de
luz es exacto ($c = 1$). La invarianza de
Lorentz~\cite{bombelli2009} no se impone externamente; emerge de
la geometr\'ia del corte.

\medskip\noindent
\texttt{diamond.rs}: \texttt{is\_causal(a, b)} implementa
literalmente~\eqref{eq:ax-lorentz}. El comentario del c\'odigo
declara: ``Light cone search: spatial radius at time offset
\texttt{dt} is \texttt{dt + margin} ($c = 1$). No fixed spatial
cutoff.''
\end{axiom}

\begin{axiom}[Densidad de Planck~\cite{planck1899,garay1995}]
\label{ax:planck-density}
El \emph{sprinkling} de Poisson opera a densidad
fundamental $\rho = 1$ evento por 4-volumen de Planck. El
di\'ametro del diamante causal es:
\begin{equation}
  T = (24\,N)^{1/4},
  \qquad
  V = \frac{T^4}{24},
  \qquad
  \rho = \frac{N}{V} = 1.
  \label{eq:ax-planck}
\end{equation}

\medskip\noindent
\texttt{diamond.rs}: \texttt{let big\_t =
(24.0 * n as f64).powf(0.25);}
\end{axiom}

\begin{axiom}[El Diferencial de Cobertura]
\label{ax:covering-differential}
El grado de libertad f\'isico fundamental es el v\'inculo de
cobertura $\prec^{\!*}$: la reducci\'on transitiva del DAG. Toda
arista $(u,v)$ tal que exista un camino $u \to w \to v$ es
redundante y se elimina. El diagrama de Hasse es la descripci\'on
can\'onica del espaciotiempo.

\medskip\noindent
\texttt{diamond.rs}: \texttt{build\_hasse\_sparse()} elimina
aristas con $A^2[u,v] > 0$. \texttt{build\_hasse\_direct()} usa
reducci\'on transitiva incremental: ``edge $(u,v)$ is a Hasse
link iff no existing link-child $z$ of $u$ satisfies
$z \prec v$.''

\stoneref{El Teorema~\ref{thm:stone-kuratowski} demuestra que
$\prec^{\!*}$ es la implicaci\'on at\'omica en el \'algebra de
Heyting, y la reducci\'on transitiva es eliminaci\'on de
corte~\cite{gentzen1935}.}
\end{axiom}

\begin{axiom}[Saturaci\'on de Valencia]
\label{ax:valence-saturation}
El grado total de cualquier evento en el diagrama de Hasse est\'a
acotado:
\begin{equation}
  \deg(u) \leq \mathcal{V}_{\max} = 15.
  \label{eq:ax-valence}
\end{equation}
Un nodo que ha acumulado 15 descendientes tiene una sombra
causal completamente poblada---cualquier candidato adicional es
transitivamente redundante.

\medskip\noindent
\texttt{diamond.rs}: \texttt{const MAX\_HASSE\_DEGREE: usize = 15;}
\end{axiom}

\begin{axiom}[Localidad Causal~\cite{malament1977}]
\label{ax:causal-locality}
La b\'usqueda de v\'inculos de Hasse y la detecci\'on de
defectos topol\'ogicos operan dentro de un horizonte local de
profundidad:
\begin{equation}
  d_{\max} = 4 \;\;\text{estratos de Hasse.}
  \label{eq:ax-locality}
\end{equation}
La f\'isica es local no por postulado sino por estructura
combinatoria: un Prisma $K_{2,N}$ tiene di\'ametro 2, y la
detecci\'on de amenazas $K_5$ requiere un salto adicional en cada
direcci\'on.

\medskip\noindent
\texttt{diamond.rs}: \texttt{const MAX\_CAUSAL\_DEPTH: i16 = 4;}
--- ``Restricting the search to $\Delta t \leq 4$ cuts the
per-node inner loop with negligible link loss ($<0.1\%$).''
\end{axiom}

\begin{axiom}[Materia Bipartita]
\label{ax:bipartite-matter}
La materia es un subgrafo bipartito completo $K_{2,N}$
denominado \textbf{Prisma Causal}, compuesto por:
\begin{enumerate}
  \item Dos \textbf{polos} (origen $u$ y destino $v$) a distancia
    de Hasse $\geq 2$.
  \item $N \geq 3$ nodos \textbf{intermedios} $\{w_1,\ldots,w_N\}$,
    mutuamente desconectados (garantizado por la propiedad
    libre de tri\'angulos), tales que $u \prec^{\!*} w_i$ y
    $w_i \prec^{\!*} v$ para todo $i$.
\end{enumerate}
La masa topol\'ogica es el entero $M = N = |\text{intermedios}|$.

\medskip\noindent
\texttt{skyrmion.rs}: \texttt{struct CausalPrism \{origin,
destination, intermediates: Vec<usize>\}}.
\texttt{const MIN\_PRISM\_SHARED: usize = 3;}

\stoneref{El Prisma es \emph{Modus Ponens} con $N$ lemas intermedios
(Teorema~\ref{thm:antimatter-tollens}).}
\end{axiom}

\begin{axiom}[Amenaza de Kuratowski~\cite{kuratowski1930}]
\label{ax:kuratowski-threat}
Un nodo externo $t \notin \mathcal{P}$ que satisface:
\begin{equation}
  t \sim u,\quad t \sim v,\quad |\{w_i \in W : t \sim w_i\}|
  \geq K_{\text{threat}} = 2,
  \label{eq:ax-threat}
\end{equation}
constituye una amenaza $K_5$ y se contrae por absorci\'on en el
polo de mayor grado. Esta operaci\'on preserva la planaridad
del grafo y es el an\'alogo discreto del confinamiento fuerte.

\medskip\noindent
\texttt{skyrmion.rs}: \texttt{const PRISM\_THREAT: usize = 2;}
--- La detecci\'on verifica
\texttt{(to\_origin \&\& to\_dest) \&\& to\_inter >=
PRISM\_THREAT}.
\end{axiom}

\begin{axiom}[Finitismo Estricto]
\label{ax:strict-finitism}
Toda acumulaci\'on interna del motor utiliza aritm\'etica entera
\texttt{u64}. Cada caminante contribuye exactamente $0$ o $1$
(delta de Kronecker) por paso de medici\'on. La \'unica
divisi\'on en punto flotante (\texttt{f64}) ocurre al final:
$P(t) = \text{count} / W$.

\medskip\noindent
\texttt{spectral.rs}: ``\textbf{Strict Finitism}: all internal
accumulation uses \texttt{u64} integer arithmetic [\ldots].
The single \texttt{f64} division \texttt{count / W} occurs only
at the final return.''
\end{axiom}

% ═══════════════════════════════════════════════════════════════════════════════
\section{Leyes}
% ═══════════════════════════════════════════════════════════════════════════════

\subsection*{L1. Ley de Reducci\'on Transitiva}
\phantomsection\label{law:transitive-reduction}

La arista $(u,v)$ es un v\'inculo de Hasse si y solo si
$A^2[u,v] = 0$: no existe ning\'un camino de dos saltos. Esta
operaci\'on produce un DAG \textbf{libre de tri\'angulos} por
necesidad matem\'atica. Consecuencia: las camarillas $K_4$ no
pueden existir---y la estructura bipartita $K_{2,N}$ es la
m\'axima permitida.

\medskip\noindent
\texttt{diamond.rs}: \texttt{build\_hasse\_sparse()} calcula
$A^2$ con producto de matrices dispersas (\texttt{sprs}) y
elimina toda arista con $A^2[i,j] > 0$.
\texttt{build\_hasse\_direct()} implementa la misma l\'ogica de
forma incremental: ``edge $(u,v)$ is a link iff no existing
link-child $z$ of $u$ satisfies $z \prec v$.''

\stoneref{La reducci\'on transitiva es eliminaci\'on de corte
(Teorema~\ref{thm:stone-kuratowski}) y No-Contradicci\'on
(Teorema~\ref{thm:non-contradiction}).}

\subsection*{L2. Ley de Masa Topol\'ogica}
\phantomsection\label{law:topological-mass}

La masa de un Prisma Causal es el n\'umero entero de nodos
intermedios:
\begin{equation}
  M = N = |W| = |\{w : u \prec^{\!*} w \prec^{\!*} v\}|
  \in \mathbb{N},\quad N \geq 3.
  \label{eq:law-mass}
\end{equation}
La masa no es un par\'ametro continuo; es un n\'umero natural.
No existe ``medio intermedio''.

\medskip\noindent
\texttt{skyrmion.rs}: \texttt{mass\_gen1 = total / ps.len()}
donde \texttt{total = ps.iter().map(|p|
p.intermediates.len()).sum()}.

\subsection*{L3. Ley de Detecci\'on 2-Hop}
\phantomsection\label{law:2hop-detection}

Para cada nodo n\'ucleo $u$, recorrer el camino $u \to w \to v$
(forward-forward):
\[
  \text{belly}(u,v) = \text{children}(u) \cap \text{parents}(v).
\]
Si $|\text{belly}| \geq 3$, se compromete un Prisma
$\mathcal{P}(u, v, \text{belly})$. Se selecciona el socio
\'optimo por tama\~no de belly (codicioso).

\medskip\noindent
\texttt{skyrmion.rs}: el bucle principal recorre
\texttt{core\_nodes}, ejecuta la b\'usqueda 2-hop con
\texttt{cands}, y calcula la intersecci\'on v\'ia
\texttt{count\_isect(nb\_u, pv)} sobre CSR ordenado en $O(D)$.

\subsection*{L4. Ley de Contracci\'on $K_5$}
\phantomsection\label{law:k5-contraction}

Si un nodo externo $t$ satisface el criterio de amenaza
(Axioma~\ref{ax:kuratowski-threat}), se absorbe en el polo de
mayor grado:
\[
  \texttt{merge\_into}[t] = \arg\max\bigl(\deg(u),\deg(v)\bigr).
\]
Las cadenas de fusiones transitivas se resuelven mediante
persecuci\'on de punteros enteros:
\texttt{while merge\_into[t] != t \{ t = merge\_into[t]; \}}.

\subsection*{L5. Ley de Clasificaci\'on Generacional}
\phantomsection\label{law:generational-classification}

Los Prismas se clasifican por su \textbf{firma topol\'ogica}: un
entero empaquetado \texttt{u32} que codifica cuatro cuartiles del
\emph{bulk momentum} (out-degree $-$ in-degree) de todos los
componentes del Prisma, normalizados y desplazados a 5~bits cada
uno:
\[
  \sigma(\mathcal{P}) =
  [q_0, q_1, q_2, q_3] \;\in\; \{0,\ldots,31\}^4.
\]
La firma m\'as abundante define Generaci\'on~1 (electr\'on), la
segunda m\'as abundante Generaci\'on~2 (mu\'on), la tercera
Generaci\'on~3 (tau).

\medskip\noindent
\texttt{skyrmion.rs}: \texttt{prism\_signature()} y el
\texttt{sig\_map: HashMap<u32, Vec<usize>>}.

\subsection*{L6. Ley de Conjugaci\'on CPT}
\phantomsection\label{law:cpt-conjugation}

La anti-firma se obtiene revirtiendo el orden e invirtiendo el
signo:
\begin{equation}
  \bar{\sigma}_i = -\sigma_{3-i} + 16.
  \label{eq:law-cpt}
\end{equation}
La Anti-Generaci\'on~1 es el conjugado CPT de la
Generaci\'on~1.

\medskip\noindent
\texttt{skyrmion.rs}: \texttt{anti\_signature(sig)} implementa
literalmente~\eqref{eq:law-cpt}.

\stoneref{La conjugaci\'on CPT es la contrapositiva l\'ogica
(Teorema~\ref{thm:antimatter-tollens}).}

\subsection*{L7. Ley del Sondeo Espectral}
\phantomsection\label{law:spectral-probe}

La geometr\'ia del grafo se mide mediante caminantes aleatorios
sobre el diagrama de Hasse no dirigido. Un caminante perezoso
(\emph{lazy walk}) permanece con probabilidad $1/2$ o salta a un
vecino uniforme con probabilidad $1/2$:
\begin{equation}
  \Pr[u_{t+1} = v \mid u_t = u] =
  \begin{cases}
    \frac{1}{2} & \text{si } v = u,\\[4pt]
    \frac{1}{2\deg(u)} & \text{si } (u,v) \in E.
  \end{cases}
  \label{eq:law-lazy}
\end{equation}
Se mide $P(t) = W^{-1}\sum_i \delta_{u_t^{(i)}, u_0^{(i)}}$ y
$\dS(t) = -2\,d(\ln P)/d(\ln t)$.

\medskip\noindent
\texttt{spectral.rs}: \texttt{run\_walkers()} con
\texttt{rng.gen\_bool(0.5)}.

\subsection*{L8. Ley del Flujo Causal (Electromagnetismo Discreto)}
\phantomsection\label{law:causal-flux}

Caminantes \emph{dirigidos} (solo causa $\to$ efecto) miden la
transmisi\'on entre firmas topol\'ogicas:
\begin{itemize}
  \item \textbf{Atracci\'on}: Gen1 $\to$ AntiGen1 (cargas opuestas).
  \item \textbf{Repulsi\'on}: Gen1 $\to$ Gen1 (cargas iguales).
\end{itemize}
La asimetr\'ia entre ambas probabilidades de transmisi\'on es
el an\'alogo discreto de la fuerza electromagn\'etica.

\medskip\noindent
\texttt{spectral.rs}: \texttt{run\_transmission\_walkers()}
con movimiento obligatorio (\texttt{pos = adj\_data[start +
rng.gen\_range(0..len)]}) y b\'usqueda binaria sobre los
conjuntos objetivo.

% ═══════════════════════════════════════════════════════════════════════════════
\section{Teoremas}
% ═══════════════════════════════════════════════════════════════════════════════

\begin{theorem}[Libertad de Tri\'angulos]
\label{thm:app-triangle-free}
La reducci\'on transitiva de cualquier DAG es libre de
tri\'angulos. Si existen aristas $u \to w$, $w \to v$, y $u \to v$
simult\'aneamente, entonces $u \to v$ es transitivamente redundante
y se elimina. Por tanto, los cliques $K_4$ no pueden existir
en el diagrama de Hasse, y $K_{2,N}$ (bipartito) es la subestructura
m\'axima permitida.

\medskip\noindent
\texttt{skyrmion.rs}, l\'inea~35: ``Therefore $K_4$ cliques
(which require triangles) CANNOT exist---zero $K_4$ is not a bug,
it is a theorem.''
\end{theorem}

\begin{proof}
Sea $H$ el diagrama de Hasse (reducci\'on transitiva) de un DAG
$G = (V, \prec)$. Supongamos, por contradicci\'on, que existen
$u, w, v$ con $u \prec^{\!*} w$, $w \prec^{\!*} v$, y
$u \prec^{\!*} v$ en $H$. Pero $u \prec w \prec v$ implica que
la arista $u \to v$ es transitivamente redundante (existe un
camino de 2~saltos $u \to w \to v$), lo cual contradice su
pertenencia a la reducci\'on transitiva. Por tanto, $H$ no
contiene tri\'angulos. Puesto que $K_4$ requiere al menos un
tri\'angulo, $K_4 \not\subset H$.
\end{proof}

\smallskip\noindent
{\small\itshape Nota: este resultado es equivalente al Principio de
No-Contradicci\'on Topol\'ogica
(Teorema~\ref{thm:non-contradiction},
Ap\'endice~\ref{app:stone-duality}), que proporciona la lectura
l\'ogica dual: la ausencia de tri\'angulos garantiza la consistencia
del sistema deductivo.}

\begin{theorem}[Completitud del 2-Hop]
\label{thm:app-2hop}
La b\'usqueda forward-forward $u \to w \to v$ encuentra
\textbf{todos} los Prismas Causales $K_{2,N}$ del n\'ucleo.
\end{theorem}

\begin{proof}
Sean $u$ y $v$ polos de un Prisma con belly
$\{w_1,\ldots,w_N\}$. Por definici\'on, $u \prec^{\!*} w_i$ y
$w_i \prec^{\!*} v$ para todo $i$. El camino
$u \to w_i \to v$ garantiza que $v$ es descubierto como candidato
2-hop a trav\'es de cualquier $w_i$. La intersecci\'on
$\text{children}(u) \cap \text{parents}(v)$ recupera el belly
completo.

\medskip\noindent
\texttt{skyrmion.rs}: ``if $u$ and $v$ are poles of a timelike
$K_{2,N}$, then every belly node $w_i$ satisfies
$u \to w_i \to v$. The path guarantees $v$ is discovered---it is
exact.''
\end{proof}

\begin{theorem}[Localidad Computacional $O(N)$]
\label{thm:app-on}
La detecci\'on de Prismas Causales escala linealmente con el
n\'umero de eventos:
\begin{equation}
  T_{\text{detect}}
  = O\bigl(|C| \times \mathcal{V}_{\max}^2\bigr)
  = O(N),
  \label{eq:thm-on}
\end{equation}
donde $|C| = N/10$ es el n\'ucleo (top $10\%$ por grado) y
$\mathcal{V}_{\max} = 15$ es la cota de valencia.
\end{theorem}

\begin{proof}
Para cada nodo n\'ucleo $u$, la colecci\'on de candidatos 2-hop
ejecuta $O(D^2)$ operaciones donde $D = \deg(u) \leq 15$. La
verificaci\'on de cada candidato (intersecci\'on de listas
ordenadas) es $O(D)$. El n\'umero de candidatos por nodo es
$O(D^2) \leq 225$. As\'i, el trabajo total por nodo es
$O(D^2 \times D) = O(D^3)$, que es $O(1)$ puesto que
$D \leq 15$. Sumando sobre $|C|$ nodos:
$T = O(|C|) = O(N/10) = O(N)$.

\medskip\noindent
\texttt{skyrmion.rs}: ``$O(N)$ scaling IS the physics.
An $O(N^2)$ algorithm would imply non-local interactions---a
universe that could not scale.''
\end{proof}

\begin{theorem}[Resoluci\'on Causal (\ref{thm:causal-resolution})]
\label{thm:app-resolution}
El n\'umero de caminantes necesario para resolver la geometr\'ia
hasta profundidad $t_{\max}$ con tolerancia~$\epsilon$ es:
\begin{equation}
  W \geq \frac{t_{\max}^{\,\dS/2}}{c\;\epsilon^2}.
\end{equation}
Para $\dS = 4$: $W \propto t_{\max}^2$ (costo cuadr\'atico).
V\'ease la demostraci\'on completa en la
Secci\'on~\ref{thm:causal-resolution}.

\medskip\noindent
\texttt{main.rs}: el n\'umero de caminantes se fija como
$W \sim N^2$ para garantizar se\~nal-ruido constante en todo el
rango de difusi\'on.
\end{theorem}

\begin{theorem}[Dimensi\'on Emergente $\dS \to 4$]
\label{thm:app-emergent-dim}
En el l\'imite de alta densidad ($N \gg 1$), la dimensi\'on
espectral del conjunto causal con Prismas converge a
$\dS = 4$ a tiempos de difusi\'on intermedios.
\end{theorem}

\begin{proof}[Esbozo]
Para un \emph{sprinkling} de Poisson en un espaciotiempo
4-dimensional, la probabilidad de retorno satisface
$P(t) \sim t^{-2}$~\cite{eichhorn2014}, de donde $\dS = -2
\times (-2) = 4$. Los Prismas Causales estabilizan $\dS$
aumentando $P(t)$ localmente (el caminante queda ``atrapado''
en los intermedios), reduciendo las fluctuaciones de $\dS$
respecto al vac\'io puro. La convergencia de $\langle \dS \rangle$
al valor $4$ se verifica num\'ericamente en la simulaci\'on
mediante promedio sobre $M$ realizaciones independientes.
\end{proof}

\begin{theorem}[Volumen de Poisson (\ref{thm:volume-convergence})]
\label{thm:app-poisson-vol}
$\mathbb{E}[N] = \rho \cdot \mathrm{Vol}(\Diamond)$,
$\mathrm{Var}[N] = \mathbb{E}[N]$, $\sigma/\mu = 1/\sqrt{N}
\to 0$. La integral topol\'ogica converge a la integral de
volumen de la Relatividad General~\cite{bombelli1987}.
\end{theorem}

\begin{theorem}[Myrheim--Meyer (\ref{thm:myrheim-meyer})]
\label{thm:app-mm}
La longitud de la cadena m\'axima
$\ell(A,B)/\rho^{1/d} \to m_d\,\tau(A,B)$ cuando
$\rho \to \infty$~\cite{myrheim1978,brightwell1991}. El conteo
de v\'inculos de cobertura converge al tiempo propio
geod\'esico.
\end{theorem}

% ═══════════════════════════════════════════════════════════════════════════════
\section{Conjeturas}
% ═══════════════════════════════════════════════════════════════════════════════

\subsection*{C1. Cuantizaci\'on del Espectro de Masas}
\phantomsection\label{conj:mass-spectrum}

El espectro de masas topol\'ogicas $\{N = 3, 4, 5, \ldots\}$
corresponde a las generaciones lept\'onicas $\{e, \mu, \tau\}$.
Las relaciones de masa $m_\mu/m_e \approx N_2/N_1$ y
$m_\tau/m_e \approx N_3/N_1$ deber\'ian coincidir con las
observaciones dentro de correcciones $O(1)$ provenientes de la
geometr\'ia del \emph{sprinkling}.

\medskip\noindent
\textit{Estado}: el histograma de masas
(\texttt{mass\_spectrum.csv}) muestra picos discretos a
$N = 3, 4, 5$. La correspondencia cuantitativa con las masas
experimentales requiere un c\'alculo de factores de forma que
a\'un no se ha implementado.

\subsection*{C2. Gravedad como Entrop\'ia de Conectividad}
\phantomsection\label{conj:gravity-entropy}

Las ecuaciones de Einstein $G_{\mu\nu} = 8\pi G\,T_{\mu\nu}$
emergen porque la densidad de Prismas modifica la probabilidad de
retorno $P(t,u)$, y v\'ia la expansi\'on del n\'ucleo del
calor~\cite{minakshisundaram1949}, esta modificaci\'on se mapea a
la curvatura de Ricci. La constante de Newton $G$ se identifica
con una combinaci\'on de la densidad fundamental~$\rho$ y el
quantum de inercia topol\'ogica~$\kappa$.

\medskip\noindent
\textit{Estado}: la Secci\'on~6.2 bosqueja la demostraci\'on.
Un esbozo de demostraci\'on se presenta en el
Teorema~\ref{thm:gravity-inference-entropy}
(Ap\'endice~\ref{app:stone-duality}, Nugget~IV) v\'ia la
entrop\'ia inferencial. La
validaci\'on num\'erica requiere extraer $R_{\mu\nu}$ de $P(t,u)$
orden por orden en la expansi\'on del calor, lo cual exige
$N > 10^7$ eventos para obtener la precisi\'on suficiente.

\subsection*{C3. Dualidad Electromagn\'etica}
\phantomsection\label{conj:em-duality}

Las probabilidades de transmisi\'on
\texttt{flux\_attraction} (Gen1 $\to$ AntiGen1) y
\texttt{flux\_repulsion} (Gen1 $\to$ Gen1) reproducen la ley de
Coulomb en el l\'imite continuo. La asimetr\'ia entre atracci\'on
y repulsi\'on es una propiedad topol\'ogica de la conjugaci\'on
CPT: firmas conjugadas facilitan la propagaci\'on causal
(atracci\'on), firmas id\'enticas la obstaculizan (repulsi\'on).

\medskip\noindent
\textit{Estado}: las simulaciones a $N = 2 \times 10^7$ muestran
consistentemente
$\texttt{flux\_repulsion} \gg \texttt{flux\_attraction}$ (ratio
$\approx 279$ a $t=1$). La asimetr\'ia es dominada por el factor
combinatorio $|\text{Gen1}| \gg |\text{AntiGen1}|$: hay $\sim 15$
veces m\'as nodos diana de repulsi\'on que de atracci\'on.
Para aislar el acoplamiento \emph{por nodo}, la normalizaci\'on
$F = \texttt{flux} / |\text{targets}|$ es necesaria; aun as\'i,
la repulsi\'on por nodo es $\sim 19\times$ mayor que la
atracci\'on, sugiriendo que la secci\'on eficaz de dispersi\'on
same-charge domina sobre la de opposite-charge en la topolog\'ia
discreta. El escalado como $1/r^2$ a\'un no se ha verificado.

\subsection*{C4. Confinamiento Fuerte como Planaridad}
\phantomsection\label{conj:strong-confinement}

La contracci\'on de amenazas $K_5$
(Axioma~\ref{ax:kuratowski-threat}) es el an\'alogo discreto del
confinamiento de color en QCD. La topolog\'ia proh\'ibe que un
nodo externo ``rompa'' la estructura bipartita del Prisma, del
mismo modo que QCD proh\'ibe aislar un quark. El \'area del tubo
de flujo conf\'inante debe escalar como la distancia entre los
polos en la m\'etrica emergente.

\medskip\noindent
\textit{Estado}: el conteo de amenazas absorbidas
(\texttt{merge\_count}) es consistente con la predicci\'on, pero
la demostraci\'on anal\'itica de la proporcionalidad \'area $\sim$
distancia permanece abierta.

\subsection*{C5. Umbral de Decoherencia}
\phantomsection\label{conj:decoherence-threshold}

Cuando el n\'umero de caminos causales internos de un subsistema
$\mathcal{S}$ viola la cota $W_{\mathcal{S}} < t^2 /
(c\,\epsilon_{\text{coh}}^2)$, la dimensi\'on espectral local
deja de converger a~4 y la descripci\'on continua se desintegra.
La decoherencia cu\'antica~\cite{zurek2003} es la p\'erdida de
resoluci\'on causal (Teorema~\ref{thm:decoherence-resolution}).

\medskip\noindent
\textit{Estado}: \textbf{demostrado} en el Teorema~\ref{thm:decoherence-resolution} (Ep\'ilogo).
Se retiene aqu\'i como referencia hist\'orica.

\subsection*{C6. Materia Oscura como Defectos Est\'eriles}
\phantomsection\label{conj:dark-matter}

El histograma de masas puede contener picos m\'as all\'a de
$N = 5$ (Generaci\'on~3) correspondientes a defectos topol\'ogicos
estables que interact\'uan gravitacionalmente (modifican $P(t,u)$)
pero no electromagn\'eticamente (no se acoplan al flujo causal
dirigido, por carecer de firma CPT alineada con las generaciones
visibles).

\medskip\noindent
\textit{Estado}: a $N > 10^6$ eventos, se observan prisms con
$|W| > 5$, pero su acoplamiento al flujo causal no se ha medido
sistem\'aticamente.

% ═══════════════════════════════════════════════════════════════════════════════
\section{Correspondencia Arquitectura $\leftrightarrow$ F\'isica}
% ═══════════════════════════════════════════════════════════════════════════════

\noindent
La siguiente tabla sintetiza c\'omo cada componente del motor de
simulaci\'on implementa una operaci\'on f\'isica fundamental:

\medskip
\begin{center}
\small
\begin{tabular}{lll}
\toprule
\textbf{M\'odulo / Funci\'on} & \textbf{Operaci\'on} &
  \textbf{F\'isica} \\
\midrule
\texttt{diamond::sprinkle()}
  & Poisson sprinkling $\rho = 1$
  & G\'enesis del espaciotiempo \\
\texttt{diamond::is\_causal()}
  & $dt^2 > |\Delta\vec{x}|^2$
  & Causalidad lorentziana \\
\texttt{diamond::build\_hasse\_*()}
  & Reducci\'on transitiva
  & $dx \to \prec^{\!*}$ \\
\texttt{skyrmion::apply\_defect()}
  & Detecci\'on $K_{2,N}$
  & Creaci\'on de materia \\
\texttt{skyrmion} (threat)
  & Contracci\'on $K_5$
  & Fuerza nuclear fuerte \\
\texttt{skyrmion::prism\_signature()}
  & Firma $[q_0,q_1,q_2,q_3]$
  & Sabor lept\'onico \\
\texttt{skyrmion::anti\_signature()}
  & $\bar{\sigma}_i = -\sigma_{3-i}$
  & Antimateria (CPT) \\
\texttt{spectral::run\_walkers()}
  & $P(t), \dS(t)$
  & Dimensi\'on del espaciotiempo \\
\texttt{spectral::run\_transmission\_*()}
  & Flujo dirigido
  & Electromagnetismo \\
\texttt{main::run\_realization()}
  & $M$ ensembles paralelos
  & Universos de Everett \\
\texttt{main::average\_ensemble()}
  & $\langle P(t)\rangle$
  & L\'imite termodin\'amico \\
\bottomrule
\end{tabular}
\end{center}

\bigskip\noindent
\begin{result}[title={Eficiencia L\'ogica como Motor de la
F\'isica}]
\phantomsection\label{res:logical-efficiency}
En la f\'isica cl\'asica, el motor que selecciona la trayectoria
f\'isica es el principio de m\'inima acci\'on: $\delta S = 0$.
En el C\'alculo de Kuratowski, este principio se reemplaza por la
\textbf{Eficiencia L\'ogica}---la evitaci\'on de contradicciones
topol\'ogicas. Toda operaci\'on del motor (\texttt{Phases 1--3})
se reduce a un \'unico imperativo: \emph{mantener la consistencia
del DAG}. La reducci\'on transitiva elimina redundancias (Ley~L1).
La detecci\'on de Prismas identifica la estructura m\'axima
compatible con la ausencia de tri\'angulos
(Teorema~\ref{thm:app-triangle-free}). La contracci\'on $K_5$
preserva la planaridad (Axioma~\ref{ax:kuratowski-threat}). El
escalado $O(N)$ demuestra que esta consistencia es mantenible
sin informaci\'on no local
(Teorema~\ref{thm:app-on}). La acci\'on cl\'asica $S[\gamma]
= -|\gamma|$ (Definici\'on~\ref{def:discrete-action}) es la
consecuencia variacional de este imperativo: la geod\'esica es la
cadena que maximiza el flujo causal sin violar la topolog\'ia. La
f\'isica no obedece leyes impuestas desde fuera; obedece la
\'unica ley que un grafo finito, causal y planar puede obedecer:
su propia consistencia l\'ogica.
\end{result}

% ███████████████████████████████████████████████████████████████████████████████
%   APPENDIX B: STONE DUALITY — PHYSICS AS LOGIC
% ███████████████████████████████████████████████████████████████████████████████

\chapter{La Dualidad de Stone: F\'isica como L\'ogica}
\label{app:stone-duality}

\begin{introduction}[Una Lente Inesperada]
  \item Este ap\'endice no a\~nade f\'isica nueva al C\'alculo de
        Kuratowski; revela una \textbf{estructura l\'ogica profunda} que
        ya estaba presente en cada axioma, ley y teorema de los
        cap\'itulos anteriores.
  \item La Dualidad de Stone~\cite{stone1936,stone1937,johnstone1982}
        establece una equivalencia categ\'orica entre topolog\'ias y
        \'algebras booleanas (l\'ogica proposicional). Si el
        espaciotiempo es una topolog\'ia (de Alexandrov sobre un DAG),
        entonces el espaciotiempo \textbf{es} un sistema l\'ogico.
  \item Presentamos cinco resultados---\emph{Golden Nuggets}---que
        emergen de leer el formalismo a trav\'es de esta dualidad.
\end{introduction}

\noindent
La correspondencia fundamental que articula este ap\'endice es:

\medskip
\begin{center}
\begin{tabular}{lll}
\toprule
\textbf{Topolog\'ia (Kuratowski)} & \textbf{L\'ogica (Stone)} &
  \textbf{F\'isica}\\
\midrule
Evento $e \in V$ (Def.~\ref{def:event})
  & Proposici\'on $\varphi$ & Punto del espaciotiempo\\
V\'inculo $\prec^{\!*}$ (Def.~\ref{def:covering})
  & Implicaci\'on $\vdash$
  & Diferencial causal $dx$\\
Diamante $\Diamond(A,B)$ (Def.~\ref{def:alexandrov})
  & Teor\'ia deductiva $\Gamma$
  & Regi\'on de integraci\'on\\
Prisma $K_{2,N}$ (Ax.~\ref{ax:bipartite-matter})
  & Lema auxiliar ($N$ pasos)
  & Part\'icula (masa $N$)\\
Amenaza $K_5$ (Ax.~\ref{ax:kuratowski-threat})
  & Paradoja l\'ogica
  & Singularidad topol\'ogica\\
Contracci\'on $K_5$ (Ley~\ref{law:k5-contraction})
  & Resoluci\'on de la paradoja
  & Fuerza nuclear fuerte\\
Cadena m\'axima $\gamma$ (Def.~\ref{def:geodesic})
  & Demostraci\'on \'optima
  & Geod\'esica (tiempo propio)\\
Anticadena $\mathcal{A}$ (Def.~\ref{def:antichain})
  & Axiomas independientes
  & Hipersuperficie de Cauchy\\
\bottomrule
\end{tabular}
\end{center}

% ═══════════════════════════════════════════════════════════════════════════════
\section{Nugget I: El Puente de Stone--Kuratowski}
% ═══════════════════════════════════════════════════════════════════════════════

\begin{theorem}[Dualidad de Stone--Kuratowski]
\label{thm:stone-kuratowski}
Sea $\mathcal{C} = (V, \prec)$ un conjunto causal finito. La
topolog\'ia de Alexandrov~\cite{sorkin1991topology} sobre $\mathcal{C}$
(cuyos abiertos son los conjuntos \emph{past-closed}:
$x \in U \wedge y \prec x \Rightarrow y \in U$) es exactamente la
topolog\'ia espectral del \'algebra de Heyting libre generada por~$V$.
En particular:

\begin{enumerate}
  \item El v\'inculo de cobertura $u \prec^{\!*} v$ \textbf{es} la
    implicaci\'on l\'ogica $\varphi_u \vdash \varphi_v$: ``la verdad
    del evento~$u$ implica causalmente la verdad del evento~$v$.''
  \item El diamante de Alexandrov $\Diamond(A,B) = \{x : A \prec x
    \prec B\}$ \textbf{es} la teor\'ia deductiva cerrada bajo
    implicaci\'on entre las proposiciones $\varphi_A$ y~$\varphi_B$.
  \item La reducci\'on transitiva (diagrama de Hasse,
    Axioma~\ref{ax:covering-differential}) \textbf{es} la
    normalizaci\'on de demostraciones (eliminaci\'on de
    corte)~\cite{gentzen1935}: elimina toda implicaci\'on
    $\varphi_u \vdash \varphi_v$ que sea derivable por
    transitividad, conservando solo los pasos inferenciales
    at\'omicos.
\end{enumerate}
\end{theorem}

\begin{proof}
La topolog\'ia de Alexandrov sobre un poset finito $P$ asigna a cada
$x \in P$ el abierto b\'asico $\downarrow\!x = \{y : y \preceq x\}$.
La funci\'on $x \mapsto \downarrow\!x$ es un isomorfismo de orden
entre $(P, \preceq)$ y el ret\'iculo de abiertos $(\mathcal{O}(P),
\subseteq)$. Por el teorema de Stone para \'algebras de
Heyting~\cite{stone1937,vickers1989}, este ret\'iculo es el \'algebra
de Lindenbaum--Tarski de una l\'ogica intuicionista cuyas
proposiciones at\'omicas son los elementos de~$V$.

Para la parte (3), observemos que en la l\'ogica intuicionista, el
\emph{Hauptsatz} de Gentzen~\cite{gentzen1935} garantiza que toda
demostraci\'on puede normalizarse eliminando ``cortes'' (usos de
transitividad). La reducci\'on transitiva $u \prec^{\!*} v$ retiene
$u \to v$ si y solo si no existe $w$ tal que $u \prec w \prec v$: es
decir, retiene exactamente las implicaciones \emph{at\'omicas}, que
son las demostraciones ya normalizadas. La eliminaci\'on de las aristas
transitivas en el DAG es literalmente la eliminaci\'on de corte en el
c\'alculo de secuentes de la l\'ogica causal.
\end{proof}

\begin{result}[title={El Espaciotiempo es un Sistema Deductivo}]
\phantomsection\label{res:stone-deductive-system}
El C\'alculo de Kuratowski no es una \emph{analog\'ia} con la l\'ogica:
\textbf{es} l\'ogica. Cada regi\'on del espaciotiempo es una teor\'ia
deductiva cerrada bajo implicaci\'on causal. La f\'isica que medimos es
la sem\'antica (los modelos) de ese sistema formal. La Dualidad de
Stone~\cite{stone1936} garantiza que ambas descripciones son
categ\'oricamente equivalentes: conocer la topolog\'ia es conocer la
l\'ogica, y viceversa.
\end{result}

% ═══════════════════════════════════════════════════════════════════════════════
\section{Nugget II: No-Contradicci\'on = Ausencia de Tri\'angulos}
% ═══════════════════════════════════════════════════════════════════════════════

\begin{theorem}[Principio de No-Contradicci\'on Topol\'ogica]
\label{thm:non-contradiction}
Sea $H = (V, E_H)$ el diagrama de Hasse de un conjunto causal
$\mathcal{C}$. Entonces $H$ es libre de tri\'angulos dirigidos.
Esta propiedad es la realizaci\'on topol\'ogica del \textbf{Principio
de No-Contradicci\'on}: en el sistema l\'ogico dual, ninguna
proposici\'on puede ser simult\'aneamente axioma y teorema derivado
respecto de otra proposici\'on.
\end{theorem}

\begin{proof}
Supongamos por contradicci\'on que existe un tri\'angulo dirigido
$u \prec^{\!*} w \prec^{\!*} v$ con $u \prec^{\!*} v$ tambi\'en
presente en~$H$. Pero $u \prec^{\!*} v$ es un v\'inculo de cobertura,
lo que exige por definici\'on que no exista ning\'un $w$ con
$u \prec w \prec v$. La presencia de $u \prec^{\!*} w \prec^{\!*} v$
viola esta definici\'on. Contradicci\'on. Luego $H$ no contiene
tri\'angulos dirigidos. $\square$

\medskip\noindent
\textit{Lectura l\'ogica.} En el sistema dual, $u \prec^{\!*} v$
significa ``$\varphi_u$ implica $\varphi_v$ en un solo paso
inferencial (axioma de implicaci\'on).'' Si existiera un camino
$\varphi_u \vdash \varphi_w \vdash \varphi_v$ Y simult\'aneamente un
axioma directo $\varphi_u \vdash \varphi_v$, entonces el mismo
v\'inculo l\'ogico ser\'ia simult\'aneamente primitivo (axiom\'atico) y
derivado (v\'ia~$w$). Esto es una contradicci\'on l\'ogica en el
meta-nivel: un mismo juicio no puede ser a la vez supuesto y
demostrado. La reducci\'on transitiva lo proh\'ibe
estructuralmente---implementando la coherencia l\'ogica del sistema.

\smallskip\noindent
\textit{Nota:} la demostraci\'on combinatoria pura de la libertad de
tri\'angulos se establece independientemente en el
Teorema~\ref{thm:app-triangle-free} (Ap\'endice~\ref{app:axioms}).
El presente resultado a\~nade la \emph{lectura l\'ogica}: la misma
propiedad grafo-te\'orica es el Principio de No-Contradicci\'on del
sistema deductivo dual.
\end{proof}

\begin{result}[title={La Geometr\'ia del Espaciotiempo Obedece la
L\'ogica Cl\'asica}]
\phantomsection\label{res:stone-non-contradiction}
La ausencia de tri\'angulos en el diagrama de Hasse no es una
curiosidad combinatoria: es la garant\'ia de que el espaciotiempo,
entendido como sistema deductivo, es \textbf{l\'ogicamente
consistente}. La reducci\'on transitiva
(Axioma~\ref{ax:covering-differential}) cumple en el DAG causal la
misma funci\'on que el Principio de No-Contradicci\'on cumple en la
l\'ogica formal: impide que una misma relaci\'on sea
simult\'aneamente at\'omica y derivada. El motor de simulaci\'on
(\texttt{build\_hasse\_sparse}) ejecuta esta prohibici\'on como su
operaci\'on m\'as costosa ($O(N^2)$ para $N \leq 15\,000$,
$O(N^{1.5})$ para $N > 15\,000$)---la consistencia l\'ogica es el
c\'omputo m\'as caro que el universo realiza.
\end{result}

% ═══════════════════════════════════════════════════════════════════════════════
\section{Nugget III: Antimateria como \emph{Modus Tollens}}
% ═══════════════════════════════════════════════════════════════════════════════

\begin{theorem}[Dualidad Materia--Antimateria como
Contrapositiva L\'ogica]
\label{thm:antimatter-tollens}
Sea $\Pi = (u, v, \{w_1, \ldots, w_N\})$ un Prisma Causal con firma
topol\'ogica $\sigma = [q_0, q_1, q_2, q_3]$. La conjugaci\'on CPT
\begin{equation}
  \bar{\sigma}_i = -\sigma_{3-i} + 16
  \label{eq:cpt-contrapositive}
\end{equation}
(\texttt{skyrmion::anti\_signature()}) es la \textbf{contrapositiva
l\'ogica} de la cadena inferencial del Prisma. Espec\'ificamente:

\begin{enumerate}
  \item El Prisma $\Pi$ codifica la cadena de \emph{Modus Ponens}:
    \begin{equation}
      \varphi_u \;\vdash\; \varphi_{w_1}, \ldots, \varphi_{w_N}
      \;\vdash\; \varphi_v.
    \end{equation}
    ``De la premisa $u$, a trav\'es de $N$ lemas intermedios, se
    concluye~$v$.''
  \item El anti-Prisma $\bar{\Pi}$ codifica el \emph{Modus Tollens}
    (contrapositiva):
    \begin{equation}
      \neg\varphi_v \;\vdash\; \neg\varphi_{w_N}, \ldots,
      \neg\varphi_{w_1} \;\vdash\; \neg\varphi_u.
    \end{equation}
    ``Si la conclusi\'on $v$ es falsa, entonces todos los
    intermedios son falsos y la premisa $u$ es falsa.''
  \item La \textbf{aniquilaci\'on} $\Pi + \bar{\Pi} \to \text{vac\'io}$
    es el colapso tautol\'ogico: la cadena inferencial se encuentra
    con su contrapositiva, los $N$ intermedios se cancelan, y
    la energ\'ia liberada $E \propto N$ es la \textbf{resoluci\'on
    combinatoria} de un tautolog\'ia---la prueba de
    $(\varphi \Rightarrow \psi) \Leftrightarrow
    (\neg\psi \Rightarrow \neg\varphi)$.
\end{enumerate}
\end{theorem}

\begin{proof}
La contrapositiva es la operaci\'on $(\varphi \Rightarrow \psi)
\mapsto (\neg\psi \Rightarrow \neg\varphi)$, v\'alida en
l\'ogica cl\'asica e intuicionista. En el poset dual, negar una
proposici\'on equivale a tomar el complemento de su filtro abierto,
lo que invierte el orden y cambia el signo de la posici\'on en el
cuartil. La firma $\sigma = [q_0, q_1, q_2, q_3]$ codifica la
distribuci\'on estad\'istica de los intermedios en cuatro cuartiles
del bulk causal. La contrapositiva invierte el orden temporal
($i \mapsto 3 - i$) y refleja la posici\'on dentro de cada cuartil
($q \mapsto -q + 16$, donde $16$ es el m\'odulo de empaquetado
\texttt{u32}). Esto es precisamente~\eqref{eq:cpt-contrapositive}.

La aniquilaci\'on libera los $N$ intermedios porque la tautolog\'ia
$(\varphi \Rightarrow \psi) \wedge (\neg\psi \Rightarrow
\neg\varphi)$ es v\'alida sin hip\'otesis---la prueba se ``colapsa''
a cero pasos, devolviendo los $N$ v\'inculos de cobertura al vac\'io.
\end{proof}

% ═══════════════════════════════════════════════════════════════════════════════
\section{Nugget IV: Gravedad como Entrop\'ia de la Inferencia}
% ═══════════════════════════════════════════════════════════════════════════════

\begin{theorem}[Ecuaciones de Einstein como Termodin\'amica
Inferencial]
\label{thm:gravity-inference-entropy}
Sea $\nu(x)$ la densidad local de Prismas Causales en torno a un
evento~$x$, y sea $\Omega(x)$ el n\'umero de cadenas causales
(demostraciones) topol\'ogicamente distintas que conectan el pasado
con el futuro de~$x$ a trav\'es de sus Prismas. Definimos la
\textbf{entrop\'ia inferencial local}:
\begin{equation}
  S_{\text{inf}}(x) \;=\; \ln \Omega(x).
  \label{eq:inferential-entropy}
\end{equation}
Entonces, en el l\'imite continuo ($\rho \to \infty$), el gradiente
de $S_{\text{inf}}$ reproduce las ecuaciones de Einstein:
\begin{equation}
  G_{\mu\nu}(x) \;=\; 8\pi G \; T_{\mu\nu}(x)
  \;\;\longleftrightarrow\;\;
  \nabla^2 S_{\text{inf}} \;=\; 8\pi G \;\nu(x),
  \label{eq:einstein-as-entropy}
\end{equation}
donde $T_{\mu\nu}$ se identifica con la densidad de defectos
topol\'ogicos $\nu(x)$ (Conjetura~C2, Ap\'endice~A).
\end{theorem}

\begin{proof}[Esbozo]
La clave es la relaci\'on entre la expansi\'on del n\'ucleo del
calor~\cite{minakshisundaram1949,gilkey1975} y la probabilidad de
retorno del caminante aleatorio
(Definici\'on~\ref{def:spectral-dimension}):
\[
  P(t, x) \;=\; \frac{1}{(4\pi t)^{d/2}}
  \Bigl(1 - \tfrac{1}{6}R(x)\,t + O(t^2)\Bigr),
\]
donde $R(x)$ es la curvatura escalar. Una regi\'on con alta densidad
de Prismas $\nu(x)$ atrapa caminantes (los intermedios act\'uan como
\emph{lazos} locales), aumentando $P(t,x)$ respecto al vac\'io. Este
aumento corresponde a una contribuci\'on negativa al t\'ermino
$-\tfrac{1}{6}R(x)t$, es decir, a curvatura positiva. Pero el
n\'umero de maneras de atrapar un caminante es exactamente
$\Omega(x)$: cada Prisma ofrece $N$ caminos alternativos por los que
el caminante puede retornar. As\'i, $P(t,x) \propto \Omega(x)$, y
por tanto $R(x) \propto -\ln \Omega(x) = -S_{\text{inf}}(x)$.

Traducido al dual l\'ogico: una regi\'on con muchos Prismas tiene
m\'as \textbf{demostraciones posibles} para conectar premisas con
conclusiones---mayor riqueza inferencial. Esta riqueza curva el
espaciotiempo. \textbf{La gravedad es la tendencia del sistema
l\'ogico a maximizar su entrop\'ia inferencial}: los eventos migran
(en el sentido del \emph{sprinkling}) hacia regiones donde hay m\'as
maneras de probar teoremas.
\end{proof}

\begin{result}[title={La Gravedad como Segundo Principio de la
L\'ogica}]
\phantomsection\label{res:stone-gravity-logic}
As\'i como el segundo principio de la termodin\'amica expresa que los
sistemas f\'isicos evolucionan hacia el m\'aximo desorden molecular,
la gravedad expresa que el sistema l\'ogico del espaciotiempo
evoluciona hacia el m\'aximo n\'umero de demostraciones posibles. La
masa curva el espacio porque \textbf{a\~nade lemas} al sistema
deductivo: cada intermedio de un Prisma es un lema que abre nuevas
v\'ias de inferencia. Cuantos m\'as lemas, m\'as demostraciones;
cuantas m\'as demostraciones, mayor entrop\'ia inferencial; mayor
entrop\'ia inferencial, mayor curvatura. Einstein no descubri\'o una
ley f\'isica: descubri\'o el \textbf{Segundo Principio de la
L\'ogica Causal}.
\end{result}

% ═══════════════════════════════════════════════════════════════════════════════
\section{Nugget V: Agujeros Negros como Horizontes G\"odelianos}
% ═══════════════════════════════════════════════════════════════════════════════

\begin{theorem}[Horizonte de Eventos como L\'imite de
G\"odel~\cite{godel1931}]
\label{thm:godel-horizon}
Sea $\mathcal{B} \subset V$ una regi\'on del DAG causal tal que su
densidad de Prismas satisface la cota de
Bekenstein~\cite{bekenstein1981}:
\begin{equation}
  \nu(\mathcal{B}) \cdot |\mathcal{B}|
  \;>\;
  \frac{|\partial \mathcal{B}|}{4\,\lP^2},
  \label{eq:bekenstein-logical}
\end{equation}
donde $|\partial \mathcal{B}|$ es el n\'umero de v\'inculos de
cobertura que cruzan la frontera de~$\mathcal{B}$. Entonces,
$\mathcal{B}$ es un \textbf{dominio g\"odeliano}: contiene
proposiciones (eventos) cuya verdad no puede ser demostrada
(alcanzada causalmente) por ninguna cadena inferencial que se
origine fuera de~$\mathcal{B}$.
\end{theorem}

\begin{proof}[Esbozo]
La cota~\eqref{eq:bekenstein-logical} establece que el n\'umero de
``teoremas internos'' (eventos causalmente conectados dentro
de~$\mathcal{B}$, cada uno multiplicado por las $N$ v\'ias
inferenciales de sus Prismas) excede la capacidad de la frontera para
transmitir demostraciones. Los $|\partial \mathcal{B}|$ v\'inculos
fronterizos son los \'unicos canales por los que una cadena
inferencial externa puede entrar en~$\mathcal{B}$. Cuando la
informaci\'on inferencial interna supera este ancho de banda, existen
necesariamente proposiciones interiores $\varphi_x$ tales que
$\varphi_x$ es \emph{verdadera} (el evento $x$ existe en el DAG)
pero \emph{indemostrable} desde el exterior---la definici\'on exacta
de una sentencia g\"odeliana~\cite{godel1931}.

El horizonte de eventos es la superficie $\partial \mathcal{B}$ donde
el ancho de banda inferencial se satura. Fuera del horizonte, el
sistema deductivo es completo (toda proposici\'on alcanzable es
demostrable). Dentro, la riqueza inferencial excede la capacidad de
comunicaci\'on: hay verdades causales que ning\'un observador externo
puede derivar.
\end{proof}

\begin{result}[title={La Paradoja de la Informaci\'on como
Incompletitud}]
\phantomsection\label{res:stone-information-paradox}
La llamada ``paradoja de la informaci\'on'' de los agujeros
negros~\cite{hawking1975} se disuelve bajo la Dualidad de Stone. No
hay paradoja porque nunca hubo p\'erdida: las proposiciones interiores
al horizonte son \emph{verdaderas pero indemostrables} desde el
exterior---exactamente como las sentencias g\"odelianas en
aritm\'etica~\cite{godel1931}. La radiaci\'on de Hawking es el
an\'alogo de los \emph{teoremas parciales} que el sistema externo
puede inferir sobre la regi\'on inaccesible: proporciona informaci\'on
estad\'istica (entrop\'ia t\'ermica) pero no puede reconstruir las
demostraciones individuales (microestados). El l\'imite de Bekenstein
$S = A/(4\lP^2)$ cuenta exactamente el n\'umero de proposiciones
interiores que son \textbf{verdaderas-pero-indemostrables} vistas
desde la frontera: es el \'indice de incompletitud g\"odeliana de la
regi\'on.
\end{result}

\bigskip
\begin{center}
\large\itshape\color{structure}
``La Dualidad de Stone revela que el C\'alculo de Kuratowski\\
no \emph{modela} la realidad con l\'ogica:\\
la realidad \textbf{es} l\'ogica ejecut\'andose a s\'i misma.\\
Cada part\'icula es un teorema, cada fuerza una regla de inferencia,\\
cada agujero negro un horizonte de incompletitud.\\
El universo no obedece leyes---\emph{es} la ley.''
\end{center}

% ███████████████████████████████████████████████████████████████████████████████
%   BIBLIOGRAPHY
% ███████████████████████████████████████████████████████████████████████████████
\backmatter
\addcontentsline{toc}{chapter}{Bibliograf\'ia}
\bibliography{modulo_synthesis_kuratowski_calculus}

\end{document}
